%%%%%%%%%%%%%%%%%%%%%%%%%%%%%%%%%%%%%%%%%
% Их сургуулийн оюутны тезис  
% LaTeX Загвар
% Version 2.3 (25/3/16)
%
% Энэ загвар нь дараах сайтаас авсан загварын монгол хувилбар юм.
% http://www.LaTeXTemplates.com
%
% Version 2.x major modifications by:
% Vel (vel@latextemplates.com)
%
% Анхдагч загварын эх үүсвэр:
% Steve Gunn (http://users.ecs.soton.ac.uk/srg/softwaretools/document/templates/)
% Sunil Patel (http://www.sunilpatel.co.uk/thesis-template/)
%
% Загварын лиценз:
% CC BY-NC-SA 3.0 (http://creativecommons.org/licenses/by-nc-sa/3.0/)
%
%%%%%%%%%%%%%%%%%%%%%%%%%%%%%%%%%%%%%%%%%

%-------------------------------------------------------------------------------
%	PACKAGES AND OTHER DOCUMENT CONFIGURATIONS
%-------------------------------------------------------------------------------

\documentclass[
11pt, % Баримтын фонтын хэмжээ, сонголт: 10pt, 11pt, 12pt
oneside, % Хоёр талаар хэвлэж үдэхээр тохируулсан. Нэг тал бол комментыг арилга
%chapterinoneline,% Нэг мөрөнд бүлгийн дугаар, нэрийг гаргах
mongolian, % babel багцын хэлний тохиргоо
singlespacing, % Мөр хоорондын зай. Сонголтууд: singlespacing, onehalfspacing, doublespacing
%draft, % Ноорог горимд шилжихийн тулд комментыг арилга(зураг, холбоос, hboxes гарахгүй)
nolistspacing, % Хэрэв мөр хоорондын зай onehalfspacing эсвэл doublespacing бол, жагсаалтын мөр хоорондын зайг single болгохын тулд комментыг арилга
%liststotoc, % Зураг/хүснэгт/бусад жагсаалтыг гарчигт оруулахын тулд комментыг арилга
%toctotoc, % Uncomment to add the main table of contents to the table of contents
%parskip, % Параграф хооронд зай оруулахын тулд комментыг арилга
%nohyperref, % hyperref багцыг ачаалахгүй бол комментыг арилга
headsepline, % Толгой мөрийн доогуур шугам татахын тулд комментыг арилга
]{MUSTThesis} % Энэ класс файл нь баримтын бүтцийг тодорхойлно

\usepackage[utf8]{inputenc} % Олон улсын тэмдэгт оруулахад хэрэгтэй
\usepackage[T2A]{fontenc} % Олон улсын тэмдэгтийн гаралтын кодчилол

\usepackage[mongolian]{babel}

% Текстийн зэрэгцүүлэлт ба paragraph-ийн уян хатан байдлыг сайжруулах
\usepackage{microtype} % Typography сайжруулах (overfull hbox багасгах)
\setlength{\emergencystretch}{3em} % Урт мөр шахах нөөц
\tolerance=1000 % Текстийн зэрэгцүүлэлтийн уян хатан байдал
\hbadness=10000 % Underfull hbox warning-ийг бууруулах

\usepackage{titlesec}
\usepackage[table]{xcolor}
\usepackage{rotating} % эргүүлэх
\usepackage{multirow}

\usepackage{tcolorbox}
\usepackage{tabularx}
\usepackage{array}
\usepackage{booktabs}
\usepackage{multirow}
\usepackage{float} % [H] placement-д хэрэгтэй
\usepackage{adjustbox} % өргөн хүснэгтийг \textwidth-д багтаах (max width)
\setlength{\emergencystretch}{2em} % урт нэр томьёоны улмаас мөр хальхаас сэргийлэх (зөөлөн)
\tcbuselibrary{skins}

% Зурганд хоосон зай үлдээхгүй float тохиргоо
\renewcommand{\topfraction}{0.95}
\renewcommand{\bottomfraction}{0.95}
\renewcommand{\textfraction}{0.05}
\renewcommand{\floatpagefraction}{0.75}
\setcounter{topnumber}{3}
\setcounter{bottomnumber}{3}
\setcounter{totalnumber}{5}

\newcolumntype{Y}{>{\raggedleft\arraybackslash}X}
\newcolumntype{Z}{>{\centering\arraybackslash}X}
\newcolumntype{L}{>{\raggedright\arraybackslash}X}

% Хүснэгтийн мөрийн өндөр (\arraystretch=1.4)-г \mainmatter-ийн дараа зөөсөн:
% нүүр хуудас, Товчилсон үгс зэрэг front matter-ийн tabular-ыг 1.0 (suvdaa шиг) байлгана.

\tcbset{tab2/.style={enhanced,fonttitle=\bfseries,fontupper=\scriptsize\rmfamily,
        colback=white,colframe=black,colbacktitle=white,
        coltitle=black,center title,sharp corners=all,
        boxsep=0pt,left=2pt,right=2pt,top=2pt,bottom=2pt,
        tabularx*={\setlength{\tabcolsep}{2pt}\renewcommand{\arraystretch}{1.0}}}}


\usepackage{amsmath}
\usepackage{amssymb}
\usepackage{braket} % Dirac notation: \ket, \bra, \braket

\usepackage{tikz} % зураг зурах
\usetikzlibrary{shapes,arrows,automata,positioning,calc,decorations.pathreplacing,trees,fit,shadows}
\usepackage{pgfplots}
\pgfplotsset{compat=1.18}
\usepackage{subcaption}

\usepackage{pdfpages} % Гарын үсэгтэй scan хуудсуудыг (PDF) шууд оруулах

\usepackage{pagecolor}
\definecolor{LigthBlue}{RGB}{102,255,255}
\definecolor{Beige}{RGB}{255,220,180}

\definecolor{OliveGreen}{cmyk}{0.64,0,0.95,0.40}
\definecolor{CadetBlue}{cmyk}{0.62,0.57,0.23,0}


\definecolor{lightlightgray}{gray}{0.9}
\usepackage{listings}
\usepackage{fancyvrb} % UTF-8 кирилл текст агуулсан скриптүүдийг хавсралтад оруулахад
\usepackage{algorithm}
\usepackage{algpseudocode} % псевдокод (algorithmic environment) — Хавсралт Б,В,Г-д ашиглана
\usepackage{newunicodechar}
% Box-drawing, math symbols, emoji-уудыг скриптийн VerbatimInput-д орлуулах
\newunicodechar{─}{-}
\newunicodechar{═}{=}
\newunicodechar{│}{|}
\newunicodechar{║}{|}
\newunicodechar{┌}{+}
\newunicodechar{┐}{+}
\newunicodechar{└}{+}
\newunicodechar{┘}{+}
\newunicodechar{├}{+}
\newunicodechar{┤}{+}
\newunicodechar{┬}{+}
\newunicodechar{┴}{+}
\newunicodechar{┼}{+}
\newunicodechar{✓}{\ensuremath{\checkmark}}
\newunicodechar{✗}{X}
\newunicodechar{→}{\ensuremath{\rightarrow}}
\newunicodechar{←}{\ensuremath{\leftarrow}}
\newunicodechar{↔}{\ensuremath{\leftrightarrow}}
\newunicodechar{≤}{\ensuremath{\leq}}
\newunicodechar{≥}{\ensuremath{\geq}}
\newunicodechar{≈}{\ensuremath{\approx}}
\newunicodechar{≠}{\ensuremath{\neq}}
\newunicodechar{—}{---}
\newunicodechar{–}{--}
\newunicodechar{×}{\ensuremath{\times}}
\newunicodechar{°}{\textdegree}
\newunicodechar{μ}{\ensuremath{\mu}}
\newunicodechar{Ω}{\ensuremath{\Omega}}
\newunicodechar{α}{\ensuremath{\alpha}}
\newunicodechar{β}{\ensuremath{\beta}}
\newunicodechar{σ}{\ensuremath{\sigma}}
\newunicodechar{π}{\ensuremath{\pi}}
\newunicodechar{∞}{\ensuremath{\infty}}
\newunicodechar{·}{\ensuremath{\cdot}}
\newunicodechar{💨}{*}
\newunicodechar{✨}{*}
\newunicodechar{⚡}{*}
\newunicodechar{🔑}{*}
\newunicodechar{🔒}{*}
\newunicodechar{📊}{*}
\renewcommand{\lstlistingname}{Эх код} % Програмын эх код хэвлэх
\lstset{
language={[x86masm]Assembler},                             % Code langugage  [x86masm]Assembler C VHDL Phyton 
basicstyle=\ttfamily,                   % Code font
keywordstyle=\color{OliveGreen},        % Keywords font ('*' = uppercase)
commentstyle=\color{gray},              % Comments font
numbers=left,                           % Line nums position
numberstyle=\tiny,                      % Line-numbers fonts
stepnumber=1,                           % Step between two line-numbers
numbersep=5pt,                          % How far are line-numbers from code
backgroundcolor=\color{lightlightgray}, % Choose background color
frame=none,                             % A frame around the code
tabsize=2,                              % Default tab size
captionpos=b,                           % Caption-position = bottom
breaklines=true,                        % Automatic line breaking?
breakatwhitespace=false,                % Automatic breaks only at whitespace?
showspaces=false,                       % Dont make spaces visible
showtabs=false,                         % Dont make tabls visible
escapeinside={(*@}{@*)},                % Монгол тэмдэгтийг lstlisting-д ашиглах
extendedchars=true,                     % UTF-8 тэмдэгт зөв уншуулах (Cyrillic, etc.)
inputencoding=utf8,
literate=%
  {а}{{\cyra}}1 {б}{{\cyrb}}1 {в}{{\cyrv}}1 {г}{{\cyrg}}1 {д}{{\cyrd}}1
  {е}{{\cyre}}1 {ё}{{\cyryo}}1 {ж}{{\cyrzh}}1 {з}{{\cyrz}}1 {и}{{\cyri}}1
  {й}{{\cyrishrt}}1 {к}{{\cyrk}}1 {л}{{\cyrl}}1 {м}{{\cyrm}}1 {н}{{\cyrn}}1
  {о}{{\cyro}}1 {п}{{\cyrp}}1 {р}{{\cyrr}}1 {с}{{\cyrs}}1 {т}{{\cyrt}}1
  {у}{{\cyru}}1 {ү}{{\cyryu}}1 {ф}{{\cyrf}}1 {х}{{\cyrh}}1 {ц}{{\cyrc}}1
  {ч}{{\cyrch}}1 {ш}{{\cyrsh}}1 {щ}{{\cyrshch}}1 {ъ}{{\cyrhrdsn}}1
  {ы}{{\cyrery}}1 {ь}{{\cyrsftsn}}1 {э}{{\cyrerev}}1 {ю}{{\cyryu}}1 {я}{{\cyrya}}1
  {ө}{{\cyrotld}}1 {Ө}{{\CYROTLD}}1 {Ү}{{\CYRYU}}1
  {А}{{\CYRA}}1 {Б}{{\CYRB}}1 {В}{{\CYRV}}1 {Г}{{\CYRG}}1 {Д}{{\CYRD}}1
  {Е}{{\CYRE}}1 {Ж}{{\CYRZH}}1 {З}{{\CYRZ}}1 {И}{{\CYRI}}1 {Й}{{\CYRISHRT}}1
  {К}{{\CYRK}}1 {Л}{{\CYRL}}1 {М}{{\CYRM}}1 {Н}{{\CYRN}}1 {О}{{\CYRO}}1
  {П}{{\CYRP}}1 {Р}{{\CYRR}}1 {С}{{\CYRS}}1 {Т}{{\CYRT}}1 {У}{{\CYRU}}1
  {Ф}{{\CYRF}}1 {Х}{{\CYRH}}1 {Ц}{{\CYRC}}1 {Ч}{{\CYRCH}}1 {Ш}{{\CYRSH}}1
  {Щ}{{\CYRSHCH}}1 {Ы}{{\CYRERY}}1 {Э}{{\CYREREV}}1 {Я}{{\CYRYA}}1
  {–}{{-}}1 {—}{{---}}1 {„}{{,,}}1 {“}{{``}}1
}

\usepackage[autostyle=false]{csquotes} % Ном зүйд хэлнээс хамаарсан хашилт оруулахад хэрэгтэй

% \usepackage[backend=bibtex,natbib=true,sorting=none,sortcites]{biblatex} % Ном зүйд bibtex -г ашиглах
% Өөрчилсөн мөр bumbayar bumabuma035@gmail.com
\usepackage[backend=biber, natbib=true, sorting=none, sortcites]{biblatex}
% babel mongolian-д biblatex-ийн mongolian.lbx файл байхгүй тул ном зүйн bibstring-ууд
% (and, in, pages, andothers, сарын нэр) bold болж рендерлэгддэг. English хэл рүү
% mapping хийснээр bibstring англиар, энгийн (тод биш) болно.
\DeclareLanguageMapping{mongolian}{english}
\addbibresource{references.bib} % Ном зүйн файл

% -----------------------------------------------------------------------------
% Гарчгийн дугаарлал ба агуулга
% -----------------------------------------------------------------------------
%   secnumdepth = 1 : 1, 1.1 дугаарлагдана; subsection (1.1.1) дугаарлагдахгүй
%   tocdepth    = 1 : Агуулгад chapter + section хүртэл л харагдана
\setcounter{secnumdepth}{1}
\setcounter{tocdepth}{1}

% Section гарчгийг зүүн талд шахсан, Bold; subsection нь дугаарлагдаагүй Bold
\titleformat{\section}{\Large\bfseries\raggedright}{\thesection}{1em}{}
\titleformat{\subsection}{\large\bfseries\raggedright}{}{0pt}{}
\titleformat{\subsubsection}{\normalsize\bfseries\raggedright}{}{0pt}{}

\newcommand{\authorshipname}{Зохиогчийн эрх хамгаалал}
\newcommand{\abbrevname}{Товчилсон үгс}
\newcommand{\constantsname}{Тогтмолууд}
\newcommand{\symbolsname}{Таних тэмдэгтүүд}
\newcommand{\acknowledgementname}{Талархал}
\newcommand{\abstractname}{Хураангуй}
\newcommand{\conclusionname}{Дүгнэлт}
% --- text2sql-new-ийн тайланд нэмсэн багц, тэмдэгт -----------------------------
\usepackage{longtable}
\usepackage{pifont}
\newunicodechar{−}{\ensuremath{-}}
\newunicodechar{±}{\ensuremath{\pm}}
\newunicodechar{Δ}{\ensuremath{\Delta}}
\newunicodechar{⇒}{\ensuremath{\Rightarrow}}
\newunicodechar{…}{\ldots}
\newunicodechar{’}{'}
\newunicodechar{√}{\ensuremath{\surd}}
\newcommand{\tsql}[1]{\texttt{\small #1}}
\newcommand{\todoresult}{\emph{(сервер дээр туршилт хийсний дараа \texttt{t2s report-latex} бөглөнө)}}
\lstdefinelanguage{SQLmn}{morekeywords={SELECT,FROM,WHERE,GROUP,BY,ORDER,JOIN,ON,AND,OR,AS,COUNT,SUM,AVG,BETWEEN,LIMIT,WITH,CASE,WHEN,THEN,ELSE,END,DISTINCT,HAVING,IN,NOT,NULL,IS},sensitive=false,morecomment=[l]{--},morestring=[b]'}

% Автоматаар үүсгэсэн төлөв (t2s report-latex): \resultsavailable г.м.
\newcommand{\resultsavailable}{0}
\newcommand{\testresultsavailable}{0}
\newcommand{\thesisresultsavailable}{1}
\newcommand{\generatedon}{2026-10-16 14:00}


%-------------------------------------------------------------------------------
%	THESIS INFORMATION
%   `end5896` — таны өөрөө бөглөх газар (text2sql-thesis/thesis/FILL_IN.md-тэй ижил журам)
%-------------------------------------------------------------------------------

\thesistitle{Монгол хэл дээрх асуултыг SQL болгон хөрвүүлэх системийн судалгаа ба хэрэгжүүлэлт: банк бус санхүүгийн байгууллагын өгөгдлийн сан дээрх туршилт}
\thesistitleeng{Mongolian Text-to-SQL: Research-Based Design, Implementation and Evaluation on a Non-Bank Financial Institution Database}
\thesistype{Бакалаврын төгсөлтийн ажил}
\authorshort{Б.Очир}
\authorlong{Батчулууны Очир}
\authorcode{B221960026}
\addresses{end5896} % ← оюутны и-мэйл
\phonenumber{end5896} % ← утас

\supervisor{end5896} % ← удирдагч багшийн зэрэг, нэр
\reader{end5896} % ← шүүмжлэгч
\advisorA{end5896} % ← зөвлөгч (байхгүй бол {})
\advisorB{}
\degreeind{end5896} % ← мэргэжлийн индекс
\degree{end5896} % ← тэнхим/хөтөлбөрийн албан нэр (ХОУ хөтөлбөр)
\subject{end5896}
\keywords{Text-to-SQL, монгол хэл, том хэлний загвар, RAG, fine-tuning, GRPO, agent, ББСБ, PostgreSQL}
\department{end5896} % ← тэнхим/хөтөлбөрийн албан нэр
\deptchair{end5896} % ← тэнхимийн эрхлэгч
\group{}
\faculty{\href{http://www.sict.edu.mn}{Мэдээлэл, Холбооны Технологийн Сургууль}}
\university{\href{http://www.must.edu.mn}{Шинжлэх Ухаан, Технологийн Их Сургууль}}

\hypersetup{pdftitle=\ttitle}
\hypersetup{pdfauthor=\shortname}
\hypersetup{pdfkeywords=\keywordnames}
\hypersetup{allcolors=black}

\begin{document}

\frontmatter
\pagestyle{plain}

%-------------------------------------------------------------------------------
%	TITLE PAGE
%-------------------------------------------------------------------------------

\begin{titlepage}
\pagecolor{LigthBlue}
\begin{center}

{\scshape\LARGE \univname\par}
{\scshape\Large \facname\par}\vspace{1 cm}

\begin{figure}[!htbp]
\centering
\includegraphics[scale=1.2]{Figures/MUST_logo.jpg}
\end{figure}

\vspace{1cm}
\hfill \large{\longname} \\

\vspace{1.5cm}

{\huge \bfseries \ttitle \par}\vspace{0.4cm}

\vspace{2.5cm}
\textsc{\Large {\thesisname}}\\

\vfill

\large {Улаанбаатар хот} \\

\end{center}
\end{titlepage}

%-------------------------------------------------------------------------------
%	SUBTITLE PAGE
%-------------------------------------------------------------------------------

\begin{titlepage}
\begin{center}
\pagecolor{white}
{\scshape\LARGE \univname\par}
{\scshape\Large \facname\par}\vspace{0.5cm}

\vspace{2cm}
\hfill \large{\deptname} \\

\vspace{2.5cm}

{\huge \bfseries \ttitle\par}\vspace{0.4cm}

\vspace{2cm}

\begin{minipage}[t] {0.9\textwidth}
\begin{flushleft}
\normalsize

Мэргэжлийн индекс: \degreeid \\
Мэргэжил: end5896 \\[2cm] % ← мэргэжлийн нэр (жишээ: Хиймэл оюун ухаан)

\begin{tabular}{l l}
\emph{Удирдагч:} & {\supname} \\
\emph{Зөвлөгч:} & {\advicenameA} \\
\emph{Гүйцэтгэгч:} & {\shortname} \\
\end{tabular}

\end{flushleft}
\end{minipage}

\vfill

\large {Улаанбаатар хот} \\
{\large 2026 он end5896 сар}\\ % ← хамгаалалтын сар

\end{center}
\end{titlepage}

%-------------------------------------------------------------------------------
%	WORK PLAN PAGE
%   Үзлэг, явц, шүүмжийн гарын үсэгтэй хуудсуудыг scan хийж signatures/huudas1..5.pdf нэрээр хадгалаад
%   доорх \includepdf мөрүүдийн %-ийг арилгана (text2sql-thesis/thesis/FILL_IN.md-тэй ижил).
%-------------------------------------------------------------------------------

\begin{titlepage}
\small
\noindent Батлав. \deptname ийн эрхлэгч:
\begin{flushright}
\makebox[6cm]{\dotfill} /\chairname/
\end{flushright}
Удирдагч:\makebox[6cm]{ \dotfill} /\supname/
\begin{center}
\vspace*{0.5cm}
\textbf{{\large \textsc{Дипломын төсөл гүйцэтгэх төлөвлөгөө}}}\\[0.5cm]
\end{center}
\noindent \textbf{Дипломын төслийн сэдэв:}\\
\textbf{Монгол}: «\ttitle» \\
\textbf{Англи}: «\ttitleng»\\

\noindent \textbf{Төслийн зорилго}: Монгол хэлээр (кирилл, латин галиг) асуусан асуултыг банк бус санхүүгийн байгууллагын
PostgreSQL өгөгдлийн сан дээр гүйцэтгэх SQL болгон хөрвүүлэх, нээлттэй эхийн загварт суурилсан системийг судалгааны
үр дүнд тулгуурлан бүтээж, хяналттай туршилтаар үнэлэх.\\

\noindent \textbf{Гүйцэтгэх оюутны овог нэр}:\makebox[4cm]{ } \shortname /\studentcode/ \\
\textbf{Холбоо барих утас}: \makebox[7cm]{ }\phonenum \\
\noindent
	\begin{tabularx}{\textwidth}{|>{\centering\arraybackslash}p{0.6cm}|X|>{\centering\arraybackslash}p{1.3cm}|>{\centering\arraybackslash}p{1.8cm}|}
	\hline
	\textbf{№} & \textbf{Ажлын бүлэг, хэсгийн нэр} & \textbf{Эзлэх хувь} & \textbf{Дуусах хугацаа} \\ \hline
	\multicolumn{4}{|l|}{\textbf{Бүлэг №1. Text2SQL-ийн онол ба судалгааны тойм}} \\ \hline
	1 & 1.1 Text2SQL асуудал ба хөгжил\newline
	    1.2 Том хэлний загварт суурилсан архитектур\newline
	    1.3 Хайлтаар баяжуулсан үүсгэлт (RAG)\newline
	    1.4 Fine-tuning: SFT ба бататгах сургалт\newline
	    1.5 Agentic Text2SQL\newline
	    1.6 Бага нөөцтэй хэл ба монгол хэлний сорилт\newline
	    1.7 Үнэлгээ ба аюулгүй байдал\newline
	    1.8 Бүлгийн дүгнэлт & 30\% & end5896 \\ \hline
	\multicolumn{4}{|l|}{\textbf{Бүлэг №2. Системийн зохиомж ба хэрэгжүүлэлт}} \\ \hline
	2 & 2.1 Шаардлага ба ерөнхий архитектур\newline
	    2.2 Өгөгдлийн сан ба text2sql давхарга\newline
	    2.3 Мэдлэгийн давхарга ба хайлт\newline
	    2.4 Загварын протокол ба SQL хамгаалалт\newline
	    2.5 Agentic урсгал\newline
	    2.6 Fine-tuning-ийн дамжлага\newline
	    2.7 Үйлчилгээ ба аюулгүй байдал\newline
	    2.8 Бүлгийн дүгнэлт & 30\% & end5896 \\ \hline
	\multicolumn{4}{|l|}{\textbf{Бүлэг №3. Туршилт ба үр дүн}} \\ \hline
	3 & 3.1 Үнэлгээний асуултын сан\newline
	    3.2 Үнэлгээний арга зүй\newline
	    3.3 Туршилтын төлөвлөгөө\newline
	    3.4 Системийн баталгаажуулалт\newline
	    3.5 Туршилтын үр дүн\newline
	    3.6 text2sql-thesis-тэй харьцуулалт\newline
	    3.7 Хязгаарлалт\newline
	    3.8 Бүлгийн дүгнэлт & 40\% & end5896 \\ \hline
	\multicolumn{4}{|l|}{\textbf{Ерөнхий дүгнэлт ба хавсралт}} \\ \hline
\end{tabularx}

\vspace{0.5cm}
Төлөвлөгөөг боловсруулсан оюутан: \makebox[3cm]{\dotfill} /\shortname/
\end{titlepage}

% Гарын үсэгтэй scan хуудсууд (байгаа бол %-ийг арилга):
%\includepdf[pages=-]{signatures/huudas1.pdf}
%\includepdf[pages=-]{signatures/huudas2.pdf}
%\includepdf[pages=-]{signatures/huudas3.pdf}
%\includepdf[pages=-]{signatures/huudas4.pdf}
%\includepdf[pages=-]{signatures/huudas5.pdf}

%-------------------------------------------------------------------------------
%	DECLARATION PAGE
%-------------------------------------------------------------------------------
\begin{declaration}
\addchaptertocentry{\authorshipname}

\noindent Миний бие \shortname, «\ttitle» сэдэвт энэ ажил нь минийх бөгөөд дараахыг нотолж байна. Үүнд:

\begin{itemize} 
\item Горилогч энэ ажлыг тус сургуулиас боловсролын зэрэг авахаар бүхэлд нь буюу голлон хийсэн болно.
\item Энэ ажлын аль нэг хэсгийг тус сургуульд эсвэл өөр байгууллагад боловсролын зэрэг, мэргэшил авахаар өмнө нь илгээсэн бол түүнийгээ тодорхой заасан болно.
\item Бусад хүмүүсийн хэвлүүлсэн ажлаас зөвлөгөө авсан бол түүнийгээ үндэслэсэн болно.
\item Бусад хүмүүсийн ажлаас ишлэл хийхдээ эх үүсвэрийг нь заасан болно.
\item Миний ажилд тусалсан голлох бүх эх үүсвэрт талархаж байна.
\item Ажлыг бусадтай хамтарсан бол алийг нь бусад хүмүүс хийсэн болохыг тодорхой заасан болно.
\end{itemize}
\bigskip
 
\noindent Гарын үсэг: \rule[-0.5em]{12.7em}{0.5pt}
\bigskip

\noindent Огноо: \rule[-0.5em]{15em}{0.5pt}

\end{declaration}

%-------------------------------------------------------------------------------
%	ABSTRACT PAGE
%-------------------------------------------------------------------------------

\addchaptertocentry{Хураангуй}

\begin{center}
{\scshape\Large \univname\par}
{\scshape\large \facname\par}\vspace{0.5cm}
{\huge\textbf{{Хураангуй}} \par}
\bigskip
{\Large{\ttitle} \par}
\bigskip

{\normalsize \shortname \par}
\addressname
\end{center}

\textit{\textbf{Түлхүүр үгс: \keywordnames}}
\bigskip

Санхүүгийн байгууллагын ажилтнууд өгөгдлийн сангаас мэдээлэл авахдаа SQL бичих чадвартай шинжээчээс хамаардаг.
Байгалийн хэлний асуултыг SQL болгон хөрвүүлэх (Text2SQL) системүүд англи хэлний BIRD, Spider зэрэг жишиг дээр
өндөр үр дүн үзүүлж байгаа ч монгол хэлэнд зориулсан систем, үнэлгээний сан, туршилтын нотолгоо бараг байхгүй.
Монгол хэлний онцлог --- залгамал үгзүй, кирилл ба латин галигийн холимог бичлэг, өгөгдлийн санд монголоор
хадгалагдсан утга --- нь англи хэлэнд боловсруулсан аргуудыг шууд хэрэглэхэд саад болдог.

Энэхүү ажлаар банк бус санхүүгийн байгууллагын (ББСБ) PostgreSQL өгөгдлийн сан дээр монгол асуултад хариулах
Text2SQL системийг судалгааны үр дүнд тулгуурлан бүтээж, хяналттай туршилтаар үнэлэх бүрэн дамжлагыг
хэрэгжүүлэв. Систем нь хоёр нээлттэй эхийн загварыг --- GRPO-оор сургасан Arctic-Text2SQL-R1-7B ба DB судлах алхамтай
(agentic) sqrl-9b --- тус бүрийн сургалтын протоколоор ажиллуулж, дараах давхаргуудтай:
(1) монгол/англи тайлбар, синоним, enum утга бүхий semantic model ба нэр томьёоны сан;
(2) BM25 ба BGE-M3 embedding-ийн хосолсон хайлт (кирилл ба галигийн түлхүүртэй утгын хайлт, нэр томьёо, жишээ);
(3) sqlglot-д суурилсан SQL хамгаалалт, read-only гүйцэтгэгч ба мөрийн түвшний хамгаалалт (RLS);
(4) алдааны мессежээр засах, олон хувилбарын санал хураалт, итгэлээр татгалзах механизм;
(5) гүйцэтгэлийн ба литерал шалгах tool, LLM шүүмжлэгчтэй agentic давталт;
(6) rejection sampling, SFT (LoRA), GRPO-д суурилсан fine-tuning-ийн дамжлага.

Үнэлгээнд зориулж 8 домэйн, 500 асуулт (dev 300, test 200) бүхий монгол eval set бэлтгэв. Асуултыг LLM агентууд
ноороглон бичиж, бие даасан LLM хянагчид шалгасан. Асуулт бүр кирилл, латин галиг, англи хувилбартай. Үүнд хариулах
боломжгүй (62) ба хоёрдмол утгатай (63) асуултыг тусгайлан оруулсан. Нэг хэв загвараас үүссэн асуултуудыг нэг
хуваалтад байлгасан тул dev ба test-ийн хооронд ойролцоо давхардал илрээгүй. Бүх gold SQL-ийг 9 өгөгдлийн сан
дээр гүйцэтгэж шалгасан. Туршилтын төлөвлөгөө 50 тохиргоог (A1--G6) «нэг
удаад нэг өөрчлөлт» зарчмаар EX, test-suite нарийвчлал, TrustSQL, McNemar тестээр харьцуулна. Мөн өмнөх ажлын
(text2sql-thesis) 150 кейсийг шилжүүлж, хоёр системийг ижил нөхцөлд харьцуулах орчин бүрдүүлсэн.

Үнэлгээний дамжлагын зөв байдлыг gold SQL-ийг загварын хариу болгож орлуулах (oracle) аргаар шалгахад A6, B5, G1
тохиргоонд, хоёр загвар, хоёр өгөгдлийн сан дээр EX 100\% гарсан. SQL injection-ий 28, монгол ба латин prompt
injection-ий 56 халдлагын бүгдийг систем хаасан. Загварын бодит оноог GPU сервер дээр (NVIDIA L4, A100) гаргана.
\ifnum\resultsavailable=1
Туршилтын үр дүнг Бүлэг~\ref{Chapter3}-т үзүүлэв.
\else
(Серверийн үр дүн хараахан ороогүй.)
\fi

%-------------------------------------------------------------------------------
%	ACKNOWLEDGEMENTS
%-------------------------------------------------------------------------------

\begin{acknowledgements}
\addchaptertocentry{\acknowledgementname}

Энэхүү ажлыг гүйцэтгэх явцад судалгааны чиглэлийг тодорхойлох, туршилтын арга зүйг боловсруулахад үнэтэй
зөвлөгөө өгсөн удирдагч багш \supname-д чин сэтгэлийн талархал илэрхийлье. % end5896 — зөвлөгч байвал нэрийг нэмж бич

Мөн ажлын турш дэмжиж, урам зориг өгсөн гэр бүл, найз нөхөддөө талархлаа илэрхийлье.

\end{acknowledgements}

%-------------------------------------------------------------------------------
%	ABBREVIATIONS
%-------------------------------------------------------------------------------

\begin{abbreviations}{l>{\raggedright\arraybackslash}p{11cm}}
\textbf{ББСБ} & Банк бус санхүүгийн байгууллага\addchaptertocentry{\abbrevname}\\
\textbf{LLM} & \textbf{L}arge \textbf{L}anguage \textbf{M}odel (том хэлний загвар)\\
\textbf{Text2SQL} & Байгалийн хэлний асуултыг SQL болгон хөрвүүлэх даалгавар\\
\textbf{SQL} & \textbf{S}tructured \textbf{Q}uery \textbf{L}anguage\\
\textbf{RAG} & \textbf{R}etrieval-\textbf{A}ugmented \textbf{G}eneration\\
\textbf{SFT} & \textbf{S}upervised \textbf{F}ine-\textbf{T}uning\\
\textbf{RL} & \textbf{R}einforcement \textbf{L}earning\\
\textbf{GRPO} & \textbf{G}roup \textbf{R}elative \textbf{P}olicy \textbf{O}ptimization\\
\textbf{LoRA} & \textbf{Lo}w-\textbf{R}ank \textbf{A}daptation\\
\textbf{CoT} & \textbf{C}hain-\textbf{o}f-\textbf{T}hought\\
\textbf{EX} & \textbf{Ex}ecution accuracy (гүйцэтгэлийн нарийвчлал)\\
\textbf{TS-EX} & \textbf{T}est-\textbf{S}uite execution accuracy\\
\textbf{BM25} & \textbf{B}est \textbf{M}atching 25 (түлхүүр үгийн хайлтын оноо)\\
\textbf{DDL} & \textbf{D}ata \textbf{D}efinition \textbf{L}anguage\\
\textbf{RLS} & \textbf{R}ow-\textbf{L}evel \textbf{S}ecurity\\
\textbf{PII} & \textbf{P}ersonally \textbf{I}dentifiable \textbf{I}nformation\\
\textbf{GPU} & \textbf{G}raphics \textbf{P}rocessing \textbf{U}nit\\
\textbf{vLLM} & LLM-ийн өндөр хурдны inference сервер\\
\textbf{BIRD} & \textbf{B}ig Bench for Large-scale Database Grounded Text-to-SQL Evaluation\\
\textbf{KPI} & \textbf{K}ey \textbf{P}erformance \textbf{I}ndicator\\
\textbf{JSONL} & JSON \textbf{L}ines (мөр бүр нэг JSON бичлэг)\\
\textbf{AML} & \textbf{A}nti-\textbf{M}oney \textbf{L}aundering (мөнгө угаахтай тэмцэх)\\
\textbf{CTE} & \textbf{C}ommon \textbf{T}able \textbf{E}xpression (SQL-ийн \texttt{WITH} дэд асуулга)\\
\textbf{HF} & \textbf{H}ugging \textbf{F}ace (загварын сан)\\
\textbf{KL} & \textbf{K}ullback--\textbf{L}eibler зөрүү\\
\textbf{KYC} & \textbf{K}now \textbf{Y}our \textbf{C}ustomer (харилцагчийг таних)\\
\textbf{LCS} & \textbf{L}ongest \textbf{C}ommon \textbf{S}ubsequence\\
\textbf{LR} & \textbf{L}earning \textbf{R}ate (сургалтын алхмын хэмжээ)\\
\textbf{LSH} & \textbf{L}ocality-\textbf{S}ensitive \textbf{H}ashing\\
\textbf{PAR30} & \textbf{P}ortfolio \textbf{a}t \textbf{R}isk 30 (30 хоногийн босготой хугацаа хэтэрсэн багцын хувь)\\
\textbf{SOV} & \textbf{S}ubject--\textbf{O}bject--\textbf{V}erb үгийн дараалал\\
\textbf{TRL} & \textbf{T}ransformer \textbf{R}einforcement \textbf{L}earning (Hugging Face-ийн сан)\\
\textbf{YAML} & YAML Ain't Markup Language (тохиргооны файлын формат)\\
\end{abbreviations}

%-------------------------------------------------------------------------------
%	ТАНИХ ТЭМДЭГТҮҮД
%-------------------------------------------------------------------------------

\begin{symbols}{l>{\raggedright\arraybackslash}p{10cm}l}
$q$            & Байгалийн хэлний асуулт (кирилл, латин галиг, англи, холимог)  & ---\addchaptertocentry{\symbolsname} \\
$S$            & Өгөгдлийн сангийн схем (хүснэгт, багана, гадаад түлхүүр)     & --- \\
$K$            & Мэдлэгийн контекст (тайлбар, нэр томьёо, утга, жишээ)          & --- \\
$\hat{s}, s^{\ast}$ & Таамагласан SQL ба лавлах (gold) SQL                       & --- \\
$D, D_i$       & Үндсэн өгөгдлийн сан ба $i$-р test-suite хуулбар ($i = 1, \dots, m$) & --- \\
$m$            & Test-suite хуулбарын тоо ($m = 8$)                              & --- \\
$N$            & Үнэлж буй асуултын тоо                                          & --- \\
$R(s, D)$      & $s$-ийг $D$ дээр гүйцэтгэсэн үр дүн (мөрийн multiset)          & --- \\
$\mathbf{1}[\cdot]$ & Индикатор функц                                            & --- \\
$\mathrm{EX}$  & Гүйцэтгэлийн нарийвчлал (тэгшитгэл~\ref{eq:ex})                 & \% \\
$\mathrm{TS}$  & Test-suite нарийвчлал, TS-EX (тэгшитгэл~\ref{eq:ts})            & \% \\
$\mathrm{Trust}_c$ & TrustSQL оноо, буруу хариултын торгууль $c \in \{1,5,10\}$ (тэгшитгэл~\ref{eq:trust}) & --- \\
$n_{10}, n_{01}$ & McNemar тестийн зөрсөн хосын тоо                             & --- \\
$\lambda$      & Хайлтад BM25 ба embedding-ийн оноог хослуулах жин ($\lambda = 0.5$) & --- \\
$G$            & GRPO-ийн нэг асуултын rollout-ийн тоо ($G=8$)                   & --- \\
$r_i, A_i$     & $i$-р rollout-ийн reward ба харьцангуй давуу тал (advantage)  & --- \\
$\beta$        & GRPO-ийн KL торгуулийн коэффициент                             & --- \\
$r, \alpha$    & LoRA-ийн rank ба масштабын коэффициент                         & --- \\
$d, k$         & LoRA-гаар засах жингийн матрицын хэмжээс ($W \in \mathbb{R}^{d \times k}$) & --- \\
$T$            & Декодчилолын температур ($T=0$: greedy)                        & --- \\
$p_{95}$       & Хариултын хугацааны 95-р процентиль                            & с \\
\end{symbols}


\tableofcontents
\listoffigures
\listoftables

\mainmatter
\renewcommand{\arraystretch}{1.3}
\pagestyle{thesis}
\myformat

\addchaptertocentry{Удиртгал}

\titleformat{\section}{\Large\bfseries\filcenter}{}{0pt}{}
\section*{Удиртгал}
\titleformat{\section}{\Large\bfseries\raggedright}{\thesection}{1em}{}
\markboth{Удиртгал}{Удиртгал}

\textbf{Сэдвийн үндэслэл.} Банк бус санхүүгийн байгууллагад (ББСБ) зээлийн багц, хугацаа хэтрэлт, эргэн төлөлт,
итгэлцэл, ерөнхий дэвтрийн орлого зэрэг өдөр тутмын асуултад хариулахад шаардлагатай мэдээлэл өгөгдлийн санд
хадгалагддаг. Ийм асуултад
хариулахын тулд менежер, эрсдэлийн шинжээч, салбарын ажилтан SQL бичих чадвартай мэргэжилтэнд хүсэлт гаргаж,
хариугаа хүлээдэг. Байгалийн хэлний асуултыг өгөгдлийн сан дээр гүйцэтгэх SQL болгон хөрвүүлэх (Text2SQL)
системүүд энэ саадыг арилгах зорилготой бөгөөд том хэлний загвар (LLM) гарч ирснээр хурдацтай хөгжиж байна
\cite{shi_2025_survey_llm_text2sql}. Англи хэлний BIRD жишиг \cite{li_2023_bird} дээр нээлттэй эхийн 7 тэрбум
орчим параметртэй загварууд том хэмжээний синтетик өгөгдлөөр сургах \cite{li_2025_omnisql} болон бататгах сургалт
(RL) хийх \cite{yao_2025_arctic_text2sql_r1} замаар хаалттай том загваруудтай өрсөлдөх түвшинд хүрсэн.

Гэвч эдгээр үр дүн монгол хэлэнд шууд шилжихгүй. Олон хэлт судалгаанууд англи бус асуулт дээр загварын нарийвчлал
мэдэгдэхүйц буурдгийг харуулсан \cite{dou_2023_multispider, pham_2026_multispider2}. Монгол хэлний хувьд нэмэлт
гурван бэрхшээл бий. Нэгдүгээрт, залгамал үгзүйн улмаас нэг үг олон хэлбэртэй («зээл», «зээлийн», «зээлээс»).
Хоёрдугаарт, хэрэглэгчид кирилл болон латин галигийг ээлжлэн, заримдаа нэг асуултад холин бичдэг («zeeliin hvselt»,
«зээлийн хүсэлт», «mobile app-аар ирсэн хүсэлт»), галигийн хэв маяг нь ч нэг биш (ү → v/u, ө → u/o). Гуравдугаарт,
өгөгдлийн санд утга монголоор (кириллээр) хадгалагддаг тул асуултад галигаар эсвэл өөр хэлбэрээр бичигдсэн утгыг
хадгалсан утгатай нь яг тааруулах шаардлагатай. Олон хэлт жишгүүд энэ асуудлыг ихэвчлэн тойрдог: BIRDTurk асуулт ба
схемийг орчуулсан ч өгөгдлийн утгыг хөндөөгүй \cite{aktas_2026_birdturk}, IndicDB-д өгөгдөл англиар хадгалагдсан
бөгөөд буруу орчуулсан нэр шүүлтийг эвддэгийг тэмдэглэсэн \cite{dawar_2026_indicdb}. Монгол хэлэнд зориулсан
Text2SQL-ийн жишиг, үнэлгээний сан, туршилтын нотолгоо бараг байхгүй байгаа нь энэ ажлын үндэслэл болов.

Санхүүгийн салбарт нэмэлт шаардлага тавигдана: систем хувийн мэдээлэл задруулахгүй, өгөгдлийг өөрчлөхгүй, салбарын
эрхийн хязгаарыг сахих ёстой. Мөн хариулж чадахгүй асуултад буруу тоо гаргахаас илүү «мэдэхгүй» гэж хариулах нь
дээр \cite{lee_2024_trustsql}. Өгөгдлийг гадагш гаргахгүйн тулд загвар байгууллагын дотор (on-prem) ажиллах нь
чухал \cite{beskorovainyi_2026_onprem_frontier}.

\textbf{Өмнөх ажил.} Энэ ажлын өмнө цахим худалдааны өгөгдлийн сан дээр монгол/англи асуултыг SQLite, Oracle,
ClickHouse диалект руу хөрвүүлэх agent системийг (text2sql-thesis) туршсан \cite{text2sql_thesis_2026}. Тэрхүү
туршилтаар (dev, $n = 21$, хоёр ажиллуулалт) монгол асуулт англиасаа 24--29 нэгжээр хоцорсон боловч McNemar-ын
$p = 0.07$--$0.22$ тул статистикийн хувьд батлагдаагүй. Retrieval-ийн recall хоёр хэлэнд ижил байсан тул зөрүүг утгын
санамжийн нарийвчлал (хэсэгчлэн) ба SQL үүсгэлт (голлон) гэж задалсан. Түүвэр жижиг бөгөөд зэрэгцээ ажиллуулалт
детерминист бус байсан тул энэ дүгнэлтийг илүү том түүвэр дээр шалгах шаардлагатай байв. Энэ ажил нь тэрхүү
дүгнэлтээс үргэлжилж, (1) ББСБ-ийн үйл ажиллагааг загварчилсан схемтэй (синтетик өгөгдөлтэй) PostgreSQL сан,
(2) хариулах боломжгүй ба хоёрдмол асуулт бүхий том eval set, (3) fine-tuning ба agentic давхарга, (4) хоёр системийг
ижил нөхцөлд харьцуулах орчныг нэмж байна.

\textbf{Судалгааны зорилго.} Монгол хэлээр (кирилл, латин галиг) асуусан асуултыг ББСБ-ийн PostgreSQL өгөгдлийн сан
дээр гүйцэтгэх SQL болгон хөрвүүлэх, нээлттэй эхийн загварт суурилсан, байгууллагын дотор ажиллах системийг
судалгааны үр дүнд тулгуурлан бүтээж, аль арга монгол хэлэнд хамгийн их үр дүн өгөхийг хяналттай туршилтаар
тогтооход оршино.

\textbf{Судалгааны зорилтууд.}
\begin{enumerate}
  \item LLM-д суурилсан Text2SQL-ийн архитектур, RAG, fine-tuning, agent, бага нөөцтэй хэлний аргууд, үнэлгээний
        судалгааг тоймлох (Бүлэг~\ref{Chapter1});
  \item ББСБ-ийн өгөгдлийн сан, мэдлэгийн давхарга, хайлт, SQL хамгаалалт, загварын протокол, agentic урсгал,
        fine-tuning-ийн дамжлагыг зохиож хэрэгжүүлэх (Бүлэг~\ref{Chapter2});
  \item 500 асуулттай монгол eval set бүтээж, хариулах боломжгүй ба хоёрдмол асуултыг оруулах, gold SQL бүрийг
        гүйцэтгэлээр баталгаажуулах;
  \item Туршилтын төлөвлөгөө (A1--G6) боловсруулж, үнэлгээний дамжлагын зөв байдлыг шалгах, хоёр загварыг
        харьцуулах (Бүлэг~\ref{Chapter3});
  \item Өмнөх ажлын 150 кейсийг шинэ систем рүү шилжүүлж, хоёр системийг ижил загвар, өгөгдөл, шалгуураар
        харьцуулах.
\end{enumerate}

\textbf{Судалгааны асуултууд.}
\begin{description}
  \item[RQ1.] Нээлттэй эхийн Text2SQL загварыг монгол асуултад шууд ашиглахад англи хувилбараас хэр хоцордог вэ?
              Латин галиг нэмэлт хэр бууралт үүсгэх вэ? (A бүлэг)
  \item[RQ2.] Схемийн тайлбар, утгын хайлт, нэр томьёо, жишээний аль нь монгол хэлэнд хамгийн их нэмэр өгөх вэ? (B бүлэг)
  \item[RQ3.] Монгол өгөгдлөөр fine-tune хийх нь RAG-аас илүү үр дүнтэй юу, хоёуланг хослуулбал ямар вэ? (C, D бүлэг)
  \item[RQ4.] Засвар, санал хураалт, татгалзалт, agentic давталт нь найдвартай байдлыг (TrustSQL) хэр сайжруулах вэ? (E, G бүлэг)
  \item[RQ5.] Шинэ систем өмнөх agent системээс ижил нөхцөлд илүү үр дүнтэй юу? (thesis benchmark)
\end{description}

\textbf{Шинэлэг тал.} (1) Бидний мэдэхээр ББСБ-ийн схем дээрх анхны монгол Text2SQL eval set: кирилл, латин галиг, англи,
холимог хувилбартай, хариулах боломжгүй ба хоёрдмол асуулттай. (2) Монгол утгын хайлтыг кирилл ба галигийн сул
түлхүүрээр хийх, дагавар таслах дүрэм. (3) Бичгээс үл хамаарах reward (F4) зэрэг монгол хэлэнд чиглэсэн санаа, мөн
англи pivot аргыг (F5, F6) монгол Text2SQL-д тохируулсан хувилбар. (4) Гүйцэтгэлийн литерал шалгах tool-той agentic
давталт. (5) Хоёр системийг ижил загвар, өгөгдөл,
шалгуураар харьцуулах дахин давтагдах орчин.

\textbf{Ажлын бүтэц.} Бүлэг~\ref{Chapter1}-т Text2SQL-ийн онол ба судалгааны тоймыг, Бүлэг~\ref{Chapter2}-т системийн
зохиомж ба хэрэгжүүлэлтийг, Бүлэг~\ref{Chapter3}-т eval set, үнэлгээний арга, туршилтын төлөвлөгөө ба үр дүнг
тайлбарлана. Хавсралтад сервер дээр ажиллуулах заавар, eval set-ийн жишээ, тохиргооны жагсаалтыг оруулав. Системийг
«Монгол хэл дээрх Text2SQL: судалгаанд суурилсан бүтээх заавар» (цаашид «заавар») ба түүний туршилтын төлөвлөгөөний
(\tsql{docs/guide/}) дагуу бүтээв.

% =============================================================================
% Бүлэг 1: Онолын судалгаа
% =============================================================================

\chapter{Text2SQL-ийн онол ба судалгааны тойм}
\label{Chapter1}
\pagecolor{white}

\section{Text2SQL асуудал ба хөгжил}
\label{sec:t2s-problem}

Text2SQL даалгаврыг дараах байдлаар томьёолно: байгалийн хэлний асуулт $q$, өгөгдлийн сангийн схем $S$, нэмэлт мэдлэг
$K$ (баганын тайлбар, бизнес нэр томьёо, жишээ) өгөгдсөн үед $R(\hat{s}, D) = R(s^{\ast}, D)$ байх SQL $\hat{s}$-ийг
үүсгэх. Энд $s^{\ast}$ нь лавлах (gold) SQL, $R(s, D)$ нь $s$-ийг $D$ өгөгдлийн сан дээр гүйцэтгэсэн үр дүн. Нэг
асуултад олон зөв SQL байж болох тул текстийн тэнцүү байдлаар бус, гүйцэтгэлийн үр дүнгээр үнэлдэг.

Shi нар Text2SQL-ийн хөгжлийг дүрэмд суурилсан, Seq2Seq, урьдчилан сургасан хэлний загвар (PLM), LLM гэсэн дөрвөн
үе шатаар тодорхойлсон \cite{shi_2025_survey_llm_text2sql}. Энд PLM-ийг нейрон үед нэгтгэж гурван үеэр авч үзэв:
\begin{itemize}
  \item \textbf{Дүрэмд суурилсан систем.} Тодорхой домэйнд гараар бичсэн дүрэм. Шинэ схемд шилжүүлэх боломж бага.
  \item \textbf{Нейрон encoder--decoder.} Spider жишиг \cite{yu_2018_spider} гарснаар олон хүснэгттэй, урьд үзээгүй
        схемд ерөнхийлөх даалгавар тавигдсан. Схем холбох (schema linking), бүтэцтэй декодчилол хөгжсөн.
  \item \textbf{Том хэлний загвар.} Prompt-оор (DIN-SQL, DAIL-SQL), олон агенттай урсгалаар (MAC-SQL, CHASE-SQL), эсвэл
        тусгай fine-tune хийсэн загвараар (OmniSQL, Arctic-Text2SQL-R1) шийддэг
        \cite{pourreza_2023_din_sql, gao_2024_dail_sql, wang_2025_mac_sql, pourreza_2024_chase_sql, li_2025_omnisql,
        yao_2025_arctic_text2sql_r1}.
\end{itemize}

BIRD жишиг \cite{li_2023_bird} нь бодит, «бохир» өгөгдөл, том схем, гадаад мэдлэг (evidence) шаарддаг асуултаар
Spider-ээс илүү бодит нөхцөлийг загварчилсан. BIRD дээр evidence-ийг хасахад нарийвчлал огцом буурдаг нь домэйны мэдлэг
чухал болохыг харуулдаг. Энэ нь мэдлэгийн давхаргын (§\ref{sec:rag}) үндэслэл болно.

\section{Том хэлний загварт суурилсан архитектур}
\label{sec:llm-arch}

Орчин үеийн Text2SQL систем ихэвчлэн дараах бүрэлдэхүүнтэй дамжлага байдаг (Зураг~\ref{fig:generic-pipeline}):
асуултыг хэвийн болгох, мэдлэг хайх (RAG), SQL үүсгэх, шалгаж гүйцэтгэх, алдаа гарвал засах, олон нэр дэвшигч SQL-ээс сонгох.

\begin{figure}[htbp]
\centering
\begin{tikzpicture}[node distance=0.45cm and 0.5cm, every node/.style={font=\footnotesize},
  box/.style={draw, rounded corners, align=center, minimum height=0.9cm, text width=2.0cm, fill=gray!8}]
  \node[box] (q) {Асуулт\\(монгол)};
  \node[box, right=of q] (n) {Хэвийн\\болгох};
  \node[box, right=of n] (r) {Мэдлэг\\хайх (RAG)};
  \node[box, right=of r] (g) {SQL\\үүсгэх (LLM)};
  \node[box, below=0.9cm of g] (v) {Шалгах,\\гүйцэтгэх};
  \node[box, left=of v] (f) {Засах\\(алдаа)};
  \node[box, left=of f] (s) {Санал\\хураах};
  \node[box, left=of s] (a) {Хариулт\\эсвэл татгалзах};
  \draw[->] (q) -- (n); \draw[->] (n) -- (r); \draw[->] (r) -- (g); \draw[->] (g) -- (v);
  \draw[->] (v) -- (f); \draw[->] (f.north east) -- (g.south west); \draw[->] (f) -- (s); \draw[->] (s) -- (a);
\end{tikzpicture}
\caption{LLM-д суурилсан Text2SQL-ийн ерөнхий дамжлага}
\label{fig:generic-pipeline}
\end{figure}

\textbf{Prompt-д суурилсан аргууд.} DIN-SQL асуудлыг хэсэгчилж (schema linking, ангилал, үүсгэлт, өөрийгөө засах)
шийддэг \cite{pourreza_2023_din_sql}. DAIL-SQL нь асуултын утгыг placeholder-оор далдалж, төстэй асуулт ба SQL-ийн
бүтцээр жишээ сонгох нь few-shot-ийн чанарыг сайжруулдгийг харуулсан \cite{gao_2024_dail_sql}.

\textbf{Олон нэр дэвшигч ба сонголт.} CHASE-SQL нь олон генератор, сонгогч, засагчийг хослуулах нь ганц генератор
эсвэл энгийн санал хураалтаас илүү болохыг ablation-оор харуулсан \cite{pourreza_2024_chase_sql}. Хамгийн энгийн
хэлбэр нь хэд хэдэн нэр дэвшигч SQL-ийг гүйцэтгээд үр дүнгээр нь санал хураах (self-consistency) юм
\cite{wang_2023_self_consistency}. Санал нийлэлтийн хэмжээ нь итгэлийн хэмжүүр болж,
бага итгэлтэй үед хариулахаас татгалзахад ашиглагдана \cite{lee_2024_trustsql}.

\textbf{Тусгай загварууд.} OmniSQL нь 2.5 сая синтетик (асуулт, SQL) хосоор Qwen2.5-Coder-ийг сургаж, жижиг загварыг
том ерөнхий загваруудтай өрсөлдөхүйц болгосон \cite{li_2025_omnisql, hui_2024_qwen25_coder}. XiYan-SQL нь M-Schema
хэмээх хагас бүтэцтэй (semi-structured) схемийн дүрслэлийг санал болгосон \cite{gao_2024_xiyan_sql}.

\section{Хайлтаар баяжуулсан үүсгэлт (RAG)}
\label{sec:rag}

RAG нь загварын параметрт байхгүй мэдлэгийг хайлтаар олж prompt-д нэмдэг \cite{lewis_2020_rag}. Text2SQL-д дөрвөн
төрлийн мэдлэг хайдаг:
\begin{enumerate}
  \item \textbf{Схемийн тайлбар.} Хүснэгт, баганын тайлбар, синоним, enum кодын нэр. Багана нь англи нэртэй
        (\tsql{status\_code}) ч монгол хэрэглэгч «төлөв», «байдал» гэж асууна.
  \item \textbf{Утга.} Асуултад гарсан «Улаанбаатар», «мобайл апп» зэрэг үгийг санд хадгалсан утгатай
        (\tsql{'MOBILE'}, \tsql{'Улаанбаатар'}) тааруулах. CHESS нь LSH-ээр \cite{talaei_2024_chess}, CodeS нь BM25 ба
        LCS-ээр \cite{li_2024_codes} утга хайдаг.
  \item \textbf{Бизнес нэр томьёо.} «Идэвхтэй багц», «хугацаа хэтэрсэн зээл» гэх мэт байгууллагын тодорхойлолт. BIRD-ийн
        evidence-тэй ижил үүрэгтэй \cite{li_2023_bird}.
  \item \textbf{Жишээ.} Төстэй асуулт ба SQL (few-shot) \cite{gao_2024_dail_sql}.
\end{enumerate}

\textbf{Хайлтын арга.} Түлхүүр үгийн BM25 \cite{robertson_2009_bm25} нь яг тохирох нэр, кодыг сайн олдог бол
embedding-д суурилсан хайлт нь утгын төстэй байдлыг барьдаг. Хоёр арга бие биеэ нөхдөг тул тэдгээрийг жинлэн
хослуулах (hybrid) нь практикт түгээмэл. Олон хэлт BGE-M3 загвар 100 гаруй хэлийг дэмждэг ч уг өгүүлэлд монгол хэлийг
тусад нь үнэлээгүй \cite{chen_2024_bge_m3}.

\textbf{Schema linking хэзээ хэрэгтэй вэ.} Хүчирхэг загварын хувьд бүтэн схем контекстэд багтах үед schema linking
(хамаарах хүснэгтийг шүүх) нь хэрэгтэй хүснэгтийг алдах эрсдэлээс болж ашиггүй, бүр хортой байж болно. Гэвч уг
судалгаанд жижиг загварт (Llama~3.1-8B) schema linking EX-ийг өсгөсөн \cite{maamari_2024_death_schema_linking}.
Манай схем (28 view) контекстэд багтдаг тул энэ ажилд бүтэн схемийг анхдагч болгож, schema linking-ийн нөлөөг 7--9B
загвар дээр тусад нь (B10) хэмжинэ.

\section{Fine-tuning: SFT ба бататгах сургалт}
\label{sec:finetune}

\textbf{LoRA.} Загварын жинг бүхэлд нь шинэчлэхийн оронд сонгосон жингийн матрицуудад бага rank-тай засвар $\Delta W = BA$
($B \in \mathbb{R}^{d \times r}$, $A \in \mathbb{R}^{r \times k}$, $r \ll \min(d,k)$) нэмж сургадаг \cite{hu_2022_lora}.
Энэ нь 7--9B загварыг нэг GPU дээр сургах боломж олгоно. QLoRA нь суурь жинг 4 битээр шахдаг \cite{dettmers_2023_qlora}.

\textbf{Rejection sampling SFT.} Загвараар нэг асуултад олон хариу үүсгээд, гүйцэтгэлийн үр дүн нь gold-той таарсныг
нь л сургалтын өгөгдөлд үлдээнэ. Иймээс зөвхөн gold SQL-тэй (reasoning-гүй) асуултаас загварын өөрийн хэв маягтай, reasoning-тэй зөв жишээ олддог.

\textbf{GRPO.} Group Relative Policy Optimization нь нэг асуултад $G$ хариу (rollout) үүсгэж, reward-ийг бүлгийн
дунджаар хэвийн болгосон давуу талаар бодлогыг шинэчилнэ \cite{shao_2024_deepseekmath}:
\begin{equation}
A_i = \frac{r_i - \operatorname{mean}(r_1,\dots,r_G)}{\operatorname{std}(r_1,\dots,r_G)}, \qquad i = 1,\dots,G.
\label{eq:grpo}
\end{equation}
Утгын функцийн (value) тусдаа загвар шаардлагагүй тул санах ой хэмнэнэ. Arctic-Text2SQL-R1 нь зөвхөн гүйцэтгэлийн
үр дүнд суурилсан энгийн reward-аар (1: зөв, 0.1: ажилласан ч буруу, 0: алдаа) GRPO хийсэн. Түүний 32B хувилбар
BIRD-ийн нэг загварын самбарт 71.83\%-иар тэргүүлсэн бол энэ ажилд ашигласан 7B хувилбар BIRD test дээр 68.5\% EX
авсан \cite{yao_2025_arctic_text2sql_r1}. Олон хэлт GRPO-ийн судалгаа ч бага өгөгдлөөр мэдэгдэхүйц
ахиц гаргаж болохыг харуулсан \cite{kattamuri_2025_contrastive_grpo}.

\section{Agentic Text2SQL}
\label{sec:agentic}

Agent гэдэг нь загвар үйлдэл хийж (tool дуудаж), орчноос ажиглалт авч, дараагийн алхмаа шийдэх давталт юм. ReAct нь
бодол ба үйлдлийг ээлжлүүлдэг \cite{yao_2023_react}. Reflexion нь өмнөх оролдлогын алдаанаас дүгнэлт гаргаж дахин
оролддог \cite{shinn_2023_reflexion}. Text2SQL-д түгээмэл tool-ууд:
(1) судлах асуулга гүйцэтгэх (\tsql{SELECT DISTINCT channel ...});
(2) утга хайх;
(3) гүйцэтгэлийн алдааг буцааж засуулах;
(4) үр дүнг шүүмжлэгчээр (critic) шалгуулах.

Sqrl-9b нь \tsql{<sql>} (судлах) ба \tsql{<answer>} (эцсийн) хэмээх хоёр үйлдэлтэй agentic протоколоор сургагдсан
загвар \cite{feyn_2026_sqrl}. Харин LLM-ийн өөрийгөө засах чадвар хязгаарлагдмал: гадаад дохиогүй (гүйцэтгэлийн
алдаа, tool-ийн үр дүн) self-correction нь заримдаа зөв хариултыг эвддэг \cite{huang_2024_self_correct}. Энэ нь agent-ийн
хамгаалалтын дүрмийн (§\ref{sec:agent}) үндэслэл болно.

\section{Бага нөөцтэй хэл ба монгол хэлний сорилт}
\label{sec:lowres}

MultiSpider нь Spider-ийг 6 хэл рүү орчуулж (англитай нийт 7 хэл), англи бус хэл дээр нарийвчлал 6.1 нэгжээр
буурдгийг харуулсан \cite{dou_2023_multispider}. Spider 2.0-ийг 8 хэлд өргөтгөсөн MultiSpider 2.0 (өөр судлаачдын
ажил) дээр хамгийн сайн reasoning загварууд ч 4\% EX авсан бол MultiSpider 1.0 дээр 60\% байв
\cite{pham_2026_multispider2}. Бага нөөцтэй хэлэнд дараах стратегиуд туршигдсан:
\begin{itemize}
  \item \textbf{Translate-test.} Асуултыг англи руу орчуулах (гадны орчуулагч эсвэл загвар өөрөө) эсвэл загварыг
        англиар сэтгүүлэх нь англи бус оролттой харьцуулахад ихэвчлэн сайжруулдаг
        \cite{etxaniz_2024_self_translate, huang_2023_xlt}. Харин Text2SQL-д нэр, байршил зэрэг утга орчуулгаар гажих
        эрсдэлтэй \cite{dawar_2026_indicdb}.
  \item \textbf{Translate-train.} Англи сургалтын өгөгдлийг (BIRD, Spider) зорилтот хэл рүү орчуулж fine-tune хийх.
        BIRDTurk нь залгамал, SOV бүтэцтэй турк хэлэнд машин орчуулга ба хүний шалгалтын дамжлага боловсруулсан
        \cite{aktas_2026_birdturk}.
  \item \textbf{Олон хэлт жишээ.} Хэл хоорондын few-shot \cite{shi_2022_xricl}.
  \item \textbf{Hint ба галиг.} IndicDB нь латинаар бичсэн, англитай холилдсон (Hinglish) асуулт ба hint-ийн нөлөөг хэмжсэн
        \cite{dawar_2026_indicdb}.
\end{itemize}

MM-Eval нь кирилл монгол хэл дээр загваруудын утга зүй, сэтгэн бодох чадвар сул байгааг \cite{zhang_2024_mm_eval},
MonCulture-Eval нь монгол соёлын мэдлэгт бичиг хоорондын зөрүү байгааг \cite{minggad_2026_monculture_eval} харуулсан. Text2SQL-ийн
хувьд монгол хэлний гол бэрхшээлийг Хүснэгт~\ref{tab:mn-challenges}-д нэгтгэв.

\begin{table}[H]
\centering\small
\caption{Монгол хэлний Text2SQL-ийн онцлог бэрхшээл ба энэ ажлын хариу}
\label{tab:mn-challenges}
\begin{tabularx}{\textwidth}{>{\raggedright\arraybackslash}p{3.2cm}LL}
\toprule
Бэрхшээл & Жишээ & Энэ ажлын хариу \\
\midrule
Залгамал үгзүй & зээл, зээлийн, зээлээс, зээлүүдийн & Дагавар таслах дүрэм (F2), embedding хайлт \\
Кирилл ба латин галиг & «hvselt», «huselt», «хүсэлт» & Галигийн сул түлхүүр, A3/F2/F3 туршилт \\
Монгол утга санд & \tsql{city = 'Улаанбаатар'} & Утгын индекс (B2), литерал шалгах tool (G) \\
Код ба нэр & \tsql{'MOBILE'} = «мобайл апп» & Semantic model-ийн enum нэр (B1) \\
Бизнес нэр томьёо & «идэвхтэй багц», «PAR30» & Glossary (B3), хоёрдмол үгийг тодруулах \\
Загварын англи төвтэй байдал & 7B загварын монгол ойлголт сул & Орчуулга (A4), pivot (F5/F6), fine-tune (C) \\
\bottomrule
\end{tabularx}
\end{table}

\section{Үнэлгээ ба аюулгүй байдал}
\label{sec:eval-theory}

\textbf{Гүйцэтгэлийн нарийвчлал (EX).} Таамаг ба gold-ийн гүйцэтгэлийн үр дүн тэнцүү асуултын хувь:
\begin{equation}
\mathrm{EX} = \frac{1}{N}\sum_{j=1}^{N} \mathbf{1}\!\left[R(\hat{s}_j, D) = R(s^{\ast}_j, D)\right].
\label{eq:ex}
\end{equation}
EX нь санамсаргүйгээр таарсан (жишээ нь шүүлт буруу ч тоо нь давхацсан) таамгийг зөв гэж тооцох эрсдэлтэй.
\textbf{Test-suite нарийвчлал} нь өгөгдлийг өөрчилсөн хэд хэдэн хуулбар DB дээр зэрэг таарахыг шаарддаг
\cite{zhong_2020_test_suite}:
\begin{equation}
\mathrm{TS} = \frac{1}{N}\sum_{j=1}^{N} \prod_{i=0}^{m} \mathbf{1}\!\left[R(\hat{s}_j, D_i) = R(s^{\ast}_j, D_i)\right],
\qquad D_0 = D.
\label{eq:ts}
\end{equation}

\textbf{Найдвартай байдал.} TrustSQL нь хариулах боломжгүй асуултыг оруулж, зөв хариулт ба зөв татгалзалтад $+1$,
хариулах боломжтой асуултаас татгалзахад $0$, буруу хариултад $-c$ оноо өгдөг \cite{lee_2024_trustsql}:
\begin{equation}
\mathrm{Trust}_c = \frac{1}{N}\sum_{j=1}^{N}
\begin{cases}
+1, & \text{зөв хариулсан эсвэл зөв татгалзсан},\\
0, & \text{хариулах боломжтой асуултаас татгалзсан},\\
-c, & \text{буруу хариулсан эсвэл хариулах боломжгүйд хариулсан}.
\end{cases}
\label{eq:trust}
\end{equation}
TrustSQL өөрөө $c \in \{0, 10, N\}$-ийг ашигласан бол энэ ажилд санхүүгийн хэрэглээнд буруу тоо их хохиролтой гэж
үзэж $c \in \{1, 5, 10\}$-ийг сонгов.

\textbf{Хоёрдмол утга ба gold-ийн алдаа.} AMBROSIA, AmbiQT нь нэг асуулт олон зөв тайлбартай байж болохыг харуулсан
\cite{saparina_2024_ambrosia, bhaskar_2023_ambiqt}. Түгээмэл жишгүүдийн gold SQL-ийн нэлээд хэсэг алдаатай
байдаг тул eval set-ийг бие даан хянах шаардлагатай \cite{jin_2026_annotation_errors, wretblad_2024_bird_noise}.
Алдааг NL2SQL-BUGs-ийн ангиллаар шинжилдэг \cite{liu_2025_nl2sql_bugs}. LLM шүүгч (FLEX) нь EX-ийн дутагдлыг
нөхөх туслах хэмжүүр байж болно \cite{kim_2025_flex}.

\textbf{Статистик ач холбогдол.} Хоёр тохиргоог ижил асуулт дээр харьцуулахад McNemar-ын тест нь зөвхөн зөрсөн
хосыг ($n_{10}$: эхнийх зөв, хоёр дахь нь буруу; $n_{01}$: эсрэгээр) ашигладаг \cite{mcnemar_1947}. Жижиг түүвэрт яг
биномиал хувилбарыг хэрэглэнэ.

\textbf{Аюулгүй байдал.} LLM-ээр дамжсан SQL injection (prompt-to-SQL) нь асуултад заавар нуух замаар өгөгдөл
устгах, бусдын мэдээлэл уншихад хүргэж болно \cite{pedro_2025_p2sql}. Хамгаалалт нь давхаргатай байх ёстой: зөвхөн
унших эрхтэй DB хэрэглэгч, мөрийн түвшний хамгаалалт (RLS), SQL-ийн бүтцийн шалгалт, гүйцэтгэлийн хязгаар. Үйлдвэрлэлийн
системүүд ч ганц загвар биш, олон бүрэлдэхүүнтэй байдаг: Uber QueryGPT нь workspace, тусгай агентуудтай бөгөөд
бүрэлдэхүүн бүрийг тусад нь үнэлдэг \cite{uber_querygpt}. LinkedIn SQL Bot нь мэдлэгийн граф ашиглаж, асуулгыг
EXPLAIN-ээр шалгадаг \cite{linkedin_sqlbot}.

\section{Бүлгийн дүгнэлт}
\label{sec:ch1-summary}

Судалгааны тоймоос дараах дүгнэлтүүд гарав:
(1) Нээлттэй эхийн 7B орчим загварууд синтетик өгөгдөл ба RL-ийн тусламжтайгаар англи Text2SQL-д өндөр түвшинд
    хүрсэн. Гэвч англи бус хэл дээр мэдэгдэхүйц буурдаг.
(2) Домэйны мэдлэг (BIRD-ийн evidence) сургалтгүйгээр мэдэгдэхүйц нэмэр өгдөг. Харин schema linking-ийн ашиг загварын
    чадвар ба схемийн хэмжээнээс хамаардаг.
(3) Translate-train ба GRPO нь англи бус хэлэнд амжилттай туршигдсан ч монгол хэлэнд шалгагдаагүй.
(4) Agent ба self-correction нь гадаад дохиотой (гүйцэтгэл, tool) үед л найдвартай.
(5) Үнэлгээнд EX-ээс гадна test-suite, татгалзалт, хоёрдмол утга, статистик тест шаардлагатай.
Эдгээрт тулгуурлан дараагийн бүлэгт системийн зохиомжийг тодорхойлно.

% =============================================================================
% Бүлэг 2: Системийн зохиомж ба хэрэгжүүлэлт
% =============================================================================

\chapter{Системийн зохиомж ба хэрэгжүүлэлт}
\label{Chapter2}
\pagecolor{white}

\section{Шаардлага ба ерөнхий архитектур}
\label{sec:arch}

Системд дараах шаардлагуудыг тавьсан:
\begin{enumerate}
  \item \textbf{Монгол хэл.} Кирилл, латин галиг (хоёр хэв маяг), англи, холимог асуултад хариулах.
  \item \textbf{Байгууллагын дотор ажиллах.} Нээлттэй эхийн загвар, өөрийн GPU сервер (vLLM \cite{kwon_2023_vllm}).
        Өгөгдөл гадагш гарахгүй.
  \item \textbf{Аюулгүй байдал.} Зөвхөн унших эрх, хувийн мэдээлэл задруулахгүй, салбарын эрхийн хязгаарыг сахих.
  \item \textbf{Найдвартай байдал.} Хариулж чадахгүй, хоёрдмол асуултыг илрүүлж татгалзах эсвэл тодруулах.
  \item \textbf{Дахин давтагдах туршилт.} Бүх тохиргоо нэг YAML файлд (\tsql{config/experiments.yaml}) тодорхойлогдож, run бүрийн мета (загварын
        хувилбар, prompt, eval set, DB-ийн fingerprint, кодын хувилбар) бичигдэх.
\end{enumerate}

Зураг~\ref{fig:architecture}-т системийн бүрэлдэхүүнийг үзүүлэв. Код нь Python (≥ 3.10) дээрх \tsql{mn\_text2sql}
багц (≈14 700 мөр), туршилтын тохиргоо нь YAML, сервер дээр ажиллуулах скриптүүд нь bash хэлбэртэй.

\begin{figure}[htbp]
\centering
\begin{tikzpicture}[node distance=0.35cm and 0.35cm, every node/.style={font=\scriptsize},
  box/.style={draw, rounded corners, align=center, minimum height=0.8cm, text width=2.8cm, fill=gray!8},
  db/.style={draw, cylinder, shape border rotate=90, aspect=0.15, align=center, minimum height=1.5cm,
             text width=2.6cm, fill=blue!5}]
  \node[box] (in) {Асуулт\\(кирилл / латин /\\англи)};
  \node[box, right=of in] (ic) {Intent шалгалт\\(PII, бичих, хоёрдмол;\\зөвхөн үйлчилгээнд)};
  \node[box, right=of ic] (ret) {Хайлт: тайлбар,\\утга, нэр томьёо,\\жишээ};
  \node[box, right=of ret] (gen) {SQL үүсгэх\\Arctic / sqrl (vLLM)};
  \node[box, below=1.1cm of gen] (guard) {SQL хамгаалалт\\(sqlglot, allow-list)};
  \node[box, left=of guard] (exec) {Read-only гүйцэтгэх\\(timeout, EXPLAIN)};
  \node[box, left=of exec] (agent) {Засвар, санал\\хураалт, agent\\(tool + critic)};
  \node[box, left=of agent] (out) {Хариулт + тайлбар\\эсвэл татгалзах};
  \node[db, below=0.7cm of exec] (pg) {PostgreSQL:\\t2s view, RLS};
  \node[db, above=0.6cm of ret] (kn) {Мэдлэг: semantic\\model, glossary, индекс};
  \draw[->] (in) -- (ic); \draw[->] (ic) -- (ret); \draw[->] (ret) -- (gen); \draw[->] (gen) -- (guard);
  \draw[->] (guard) -- (exec); \draw[->] (exec) -- (agent); \draw[->] (agent) -- (out);
  \draw[<->, dashed] (exec) -- (pg); \draw[<->, dashed] (ret) -- (kn);
  \draw[->] (agent.north) -- ++(0,0.5) -| node[pos=0.25, below, font=\tiny]{алдаа бол засвар} ([xshift=-0.7cm]gen.south);
\end{tikzpicture}
\caption{Системийн бүрэлдэхүүн ба өгөгдлийн урсгал (алдаа гарвал засвар нь SQL үүсгэх алхам руу буцна)}
\label{fig:architecture}
\end{figure}

\section{Өгөгдлийн сан ба text2sql давхарга}
\label{sec:db}

\textbf{Схем.} Өмнөх ажлын репозиторт (text2sql-thesis) байсан, тэр ажлын туршилтад ашиглагдаагүй ББСБ-ийн PostgreSQL
схемийг ашиглав. 13 бизнес схемд нийт 189 хүснэгт ба view бий: байгууллагын үндсэн мэдээлэл (салбар, ажилтан),
харилцагч, KYC, бүтээгдэхүүн, зээл, барьцаа, данс ба гүйлгээ, хадгаламж, барагдуулалт (collection), эрсдэл, нягтлан
бодох бүртгэл, тайлан, систем. Ачаалах үед илэрсэн алдааг зөвхөн хуулбар дээр засаж, \tsql{db/CHANGES.md}-д бүртгэсэн.
Монгол Улсын хуулиар ББСБ иргэн, хуулийн этгээдээс хадгаламж авахыг хориглодог (харин итгэлцлийн үйлчилгээ эрхэлж
болно) \cite{mn_2002_nbfi_law}. Эх схемийн хадгаламжийн (\tsql{deposit}) хэсэг энэ утгаараа бодит ББСБ-ээс зөрөх
боловч итгэлцэлтэй төстэй бүтээгдэхүүний жишээ болгон туршилтад хэвээр үлдээв.

\textbf{Синтетик өгөгдөл.} Бодит харилцагчийн мэдээлэл ашиглах боломжгүй тул Python-оор детерминист (seed 42)
синтетик өгөгдөл үүсгэв. «Өнөөдөр» нь 2026-06-30, түүх 2024-01-01-ээс эхэлнэ. Үндсэн сан: 3000 харилцагч, 2666 зээл,
212 566 гүйлгээ, 17 салбар, 158 ажилтан. Зээлийн симуляц нь хуваарь, төлөлт, хугацаа хэтрэлт, бүтцийн өөрчлөлт,
шүүхийн шат, хасалт, ерөнхий дэвтрийн бичилтийг хамарна. DB-ийн функцүүд \tsql{CURRENT\_DATE} ашигладаг тул «өнөөдөр»-ийн
утга ачаалсан өдрөөс хамаарахгүй байлгахын тулд өгөгдлийг Python-д тооцож ачаалсан.

\textbf{text2sql давхарга (\tsql{t2s} схем).} Загварт 189 хүснэгт ба view-ийг шууд үзүүлэхийн оронд 28 view бүхий
тусгай схем бүтээв. View бүр нэг бизнес объектод тохирно (\tsql{loans}, \tsql{loan\_payments}, \tsql{collection\_cases},
\tsql{deposits}, \tsql{gl\_entries} г.м.). Код бүрийн хажууд нэрийн багана нэмж (\tsql{status\_code} → \tsql{status\_name}),
олон валютын дүнг төгрөгт хөрвүүлсэн багана (\tsql{*\_mnt}) нэмж, хувийн мэдээллийн багануудыг (регистр, утас,
хаяг) хассан. Үлдсэн харилцагчийн нэрийг semantic model-д хувийн мэдээлэл гэж тэмдэглэсэн. Туслах функц нь
\tsql{t2s.business\_date()} ба \tsql{t2s.fx\_to\_mnt()}.

\textbf{Эрх.} Bot хэрэглэгч (\tsql{t2s\_bot}) зөвхөн \tsql{t2s} схемд SELECT эрхтэй. Холболт бүр read-only
транзакц, 20 секундийн statement timeout, 2 секундийн lock timeout, 1 GB түр файлын хязгаартай. Мөрийн түвшний
хамгаалалт (RLS) нь салбарын хэрэглэгчид зөвхөн өөрийн салбарын мөрийг харуулна
(\tsql{t2s\_branch\_demo} хэрэглэгчээр туршсан).

\textbf{Test-suite хуулбарууд.} EX-ийн санамсаргүй тохиролтыг шалгахын тулд өөр seed-тэй, 0.25 масштабтай 8 хуулбар
DB (\tsql{nbfi\_ts1..8}) үүсгэв. Хуулбар бүр өөр төрлийн «хавх» агуулна: NULL утгыг олшруулсан, дүнг бөөрөнхийлж тэнцүү
утга үүсгэсэн, хугацаа хэтрэлт ба огнооны хил (30/31, 90/91 хоног, сарын эхэн ба сүүл), зарим хүснэгтийг хоосолсон
(хасалт, AML, батлан даагч, бүтцийн өөрчлөлт, хугацаанаас өмнөх татан авалт), эдгээрийн хослол ба энгийн өөр seed.

\section{Мэдлэгийн давхарга ба хайлт}
\label{sec:knowledge}

\textbf{Semantic model} (\tsql{knowledge/semantic\_model.yaml}). 28 view, 425 баганын монгол ба англи тайлбар,
кирилл синоним, enum кодын монгол нэр (\tsql{MOBILE = Мобайл апп}), гадаад түлхүүр, хувийн мэдээллийн тэмдэглэгээ.
Багана, хүснэгтийн нэр англиар үлддэг. Харин тайлбар нь DDL-ийн comment хэлбэрээр prompt-д орно
(Эх код~\ref{lst:ddl}).

\begin{lstlisting}[language=SQLmn, basicstyle=\ttfamily\scriptsize, caption={B1-ийн тайлбартай DDL-ийн хэсэг (prompt-д орох хэлбэр, товчилсон)}, label={lst:ddl}]
CREATE TABLE loan_applications (
  application_id TEXT,  -- (*@Хүсэлтийн ID@*)
  application_no TEXT,  -- (*@Хүсэлтийн дугаар@*) (LA-2025-000123); e.g. LA-2024-000001, LA-2024-000002
  status_code TEXT,  -- (*@Хүсэлтийн төлөв;@*) DRAFT=(*@Ноорог@*), SUBMITTED=(*@Илгээгдсэн@*), ..., REJECTED=(*@Татгалзсан@*), ...
  channel TEXT,  -- (*@Хүсэлт ирсэн суваг;@*) BRANCH=(*@Салбар@*), MOBILE=(*@Мобайл апп@*), WEB=(*@Вэб@*), AGENT=(*@Агент@*), ...
  requested_amount NUMERIC,  -- (*@Хүссэн дүн@*)
  ...
\end{lstlisting}

\textbf{Glossary} (\tsql{knowledge/glossary.yaml}). Байгууллагын 45 нэр томьёо: тодорхойлолт, каноник SQL, хамаарах
хүснэгт. Үүнээс 7 нь хоёрдмол («идэвхтэй зээл» = идэвхтэй багц уу, эсвэл зөвхөн хэвийн төлөвтэй юу).
Үйлчилгээний горимд intent шалгалт хоёрдмол үгийг тодотголгүй асуувал тодруулах асуулт буцаана. Туршилтад энэ
шалгалт унтраалттай тул хоёрдмол асуултад загвар өөрөө хариулна (§\ref{sec:eval-method}).

\textbf{KPI-ийн semantic layer} (\tsql{knowledge/metrics.yaml}). 20 гол үзүүлэлт, 14 хэмжээс (185 зөвшөөрөгдсөн хос).
F8 туршилтад жижиг загвар нь асуултыг \{metric, хэмжээс, шүүлт\} JSON болгож, SQL-ийг кодоор compile хийнэ.

\textbf{Индекс ба хайлт.} \tsql{index-build} нь дөрвөн индекс үүсгэнэ:
(1) схемийн баримт;
(2) текст баганын давтагдашгүй утга (хувийн мэдээллийн баганаас бусад, ≤400 утгатай);
(3) нэр томьёо;
(4) 160 жишээ асуулт.
Хайлт нь BM25 ба BGE-M3 embedding-ийн оноог $\lambda = 0.5$ жингээр хослуулна. Утгын хайлт гурван горимтой:
\begin{itemize}
  \item \textbf{basic} (B2): кирилл түлхүүр ба галигийн «сул түлхүүр». Сул түлхүүр нь о/ө/у/ү-г нэг үсэг болгож,
        kh/h, ts/c, w/v-г нэгтгэж, давхар үсгийг дан болгоно. Латин «v» нь ү эсвэл в байж болох тул хоёр хувилбараар хайна;
  \item \textbf{rule} (F2): + дагавар таслах (-ийн, -аас, -д …) ба галигийн хоёр хэв маягийн хувиргалт;
  \item \textbf{learned} (F3): + хэрэглэгчийн лог ба засвараас сурсан галиг → кирилл хос.
\end{itemize}
Few-shot жишээг DAIL-SQL маягаар асуултын утгыг \tsql{<утга>}, \tsql{<тоо>}-оор далдалж төстэй байдлаар сонгоно
\cite{gao_2024_dail_sql}. Схемийн тайлбар токены төсвөөс (13 000) хэтэрвэл бүх хүснэгт DDL-д үлдэж, тайлбарыг зөвхөн
асуулттай хамгийн холбоотой хүснэгтүүдэд өгнө.

\section{Загварын протокол ба SQL хамгаалалт}
\label{sec:protocol}

Туршилтын хоёр гол загварыг \textbf{тус бүрийн сургалтын протоколоор} ажиллуулна. Учир нь RL-ээр сургасан загвар
өөр prompt хэлбэрт муудаж болох бөгөөд sqrl-9b-ийн карт ч сургалтын протоколоо яг мөрдөхийг зөвлөдөг
\cite{feyn_2026_sqrl}.

\textbf{Arctic-Text2SQL-R1-7B} (Qwen2.5-Coder-7B-Instruct-д суурилсан, OmniSQL-оос эхлүүлж GRPO хийсэн)
\cite{yao_2025_arctic_text2sql_r1, snowflake_2025_arctic_card}. A1-ийн prompt нь өгүүллийн Зураг C.1-тэй ижил (бусад
тохиргоонд жишээ ба нэмэлт зааврын хэсэг нэмэгдэнэ): system нь «data science expert», user нь
Database Engine / Database Schema / Question / Instructions / Output Format хэсэгтэй. Assistant-ийн хариуг
«Let me solve this step by step.\textbackslash n<think>» гэж эхлүүлж (prefill), SQL-ийг \tsql{<answer>} доторх
sql код блокоос авна. RAG-ийн агуулга нь DDL-ийн comment ба асуултын өмнөх evidence хэсэгт орно.

\textbf{sqrl-9b} (Qwen3.5-9B суурьтай agentic загвар) \cite{feyn_2026_sqrl}. \tsql{<schema>}, \tsql{<evidence>} хэсэгтэй
system prompt. Загвар \tsql{<sql>...</sql>}-ээр DB-г судалж (≤5 алхам), ажиглалтыг \tsql{<observation>}-оор
авна. Эцэст нь \tsql{<answer>}-ээр хариулна.

\textbf{Диалект.} Хоёр загвар голчлон SQLite дээр сургагдсан. Иймд гурван аргыг туршина:
\begin{enumerate}
  \item SQLite протоколоор шууд асууж, гарсан SQL-ийг PostgreSQL дээр гүйцэтгэх (A1);
  \item prompt-д PostgreSQL-ийг заах (A6);
  \item SQLite-аар үүсгээд sqlglot-оор PostgreSQL руу хөрвүүлэх (A7) \cite{mao_2026_sqlglot}.
\end{enumerate}

\textbf{SQL хамгаалалт.} Загварын гаргасан SQL бүр (sqrl-ийн судлах асуулга ч) дараах шалгалтыг давна:
\begin{enumerate}
  \item sqlglot-оор задлах; ганц \tsql{SELECT}/\tsql{WITH} асуулга байх;
  \item өгөгдөл өөрчлөх, DDL, \tsql{SET}, \tsql{COPY}, \tsql{SELECT INTO}, \tsql{FOR UPDATE} хориотой;
  \item хориотой функцүүд (\tsql{pg\_sleep}, \tsql{dblink}, \tsql{pg\_read\_file}, \tsql{set\_config} г.м. 35 гаруй);
  \item зөвхөн \tsql{t2s} схемийн бүртгэлтэй view (alias-тай баганыг allow-list-ээр, бусдыг EXPLAIN-ээр шалгана);
  \item \tsql{CURRENT\_DATE}, \tsql{now()}, \tsql{date('now')}, нэг аргументтай \tsql{age()}-ийг бизнесийн
        өдрөөр дахин бичих (үнэлгээ хуанлийн өдрөөс хамаарахгүй);
  \item EXPLAIN-ийн өртгийн дээд хязгаар ($5 \times 10^7$), мөрийн хязгаар.
\end{enumerate}
Энэ нь эхний шүүлтүүр бөгөөд DB-ийн эрх, RLS, read-only транзакц нь хоёр дахь, бие даасан хамгаалалт болно.

\textbf{Засвар, санал хураалт, татгалзалт.} E1 нь алдааны мессеж (эсвэл хоосон үр дүн) буцааж 2 удаа хүртэл засуулна
(төлөвлөгөөнд 1--3). Загвар SQL бичихээс татгалзсан (SQL-гүй бүрэн хариу өгсөн) бол засуулахгүй, «татгалзсан» гэж
тооцно. Тасарсан гаралт ба SQL олдоогүй форматын алдаа нь алдаа хэвээр бөгөөд E1 засуулна.
E2 нь 8 нэр дэвшигч SQL үүсгэж (sampling), гүйцэтгэлийн үр дүнгийн hash-аар санал хураана. Олонхын хувь нь итгэлийн
хэмжүүр болно. E3 нь итгэл босгоос доош бол хариулахаас татгалзана.

\section{Agentic урсгал}
\label{sec:agent}

Өмнөх ажлын Planner → Writer → Critic агентын санааг шинэ системийн хамгаалалт, хайлт дээр дахин зохиов
(Зураг~\ref{fig:agent}). Загвар өөрийгөө засах чадвар гадаад дохиогүй үед найдваргүй
\cite{huang_2024_self_correct}. Иймд agent-ийн гол дохио нь \textbf{детерминист tool}-ууд байна:
\begin{itemize}
  \item \textbf{Гүйцэтгэлийн шалгалт.} Алдаа гарвал заавал засуулах хатуу олдвор болно. Харин хоосон үр дүн ба бүгд NULL нь
        зөөлөн олдвор: зөв хариу хоосон байж болох тул шүүмжлэгчид л дамжуулна.
  \item \textbf{Литерал шалгах tool.} \tsql{WHERE}, \tsql{IN}, \tsql{LIKE} дахь текст утга тухайн баганад үнэхээр хадгалагдсан
        эсэхийг DB-ээс шалгана. Байхгүй бол тухайн баганын хамгийн ойр хадгалсан утгуудыг санал болгоно. Хувийн
        мэдээллийн баганын утгыг харуулахгүй. Үр дүн хоосон биш үед олдвор зөөлөн байна.
  \item \textbf{LLM шүүмжлэгч.} Асуулт, SQL, үр дүнгийн эхний 8 мөр, tool-ийн олдворыг өгөхөд загвар
        «VERDICT: OK» эсвэл «VERDICT: REVISE --- …» гэж дүгнэнэ. Дүгнэлтгүй хариуг SQL өөрчлөхгүйгээр
        тусад нь бүртгэнэ.
\end{itemize}
Засварласан SQL гүйцэтгэгдэхгүй, эсвэл ажиллаж байсан үр дүнг хоосон болговол түүнийг үгүйсгэж, өмнөх SQL-ийг
хадгална. Загвар SQL бичихээс татгалзсан үед (хариулах боломжгүй асуулт) agent SQL бичүүлэх гэж шахахгүй.
Сонголтоор planner нь SQL бичихээс өмнө англи төлөвлөгөө (хүснэгт, join, шүүлт, агрегат) гаргана.

\begin{figure}[htbp]
\centering
\begin{tikzpicture}[node distance=0.5cm and 0.6cm, every node/.style={font=\scriptsize},
  box/.style={draw, rounded corners, align=center, minimum height=0.8cm, text width=2.4cm, fill=gray!8},
  dec/.style={draw, diamond, aspect=2.4, align=center, inner sep=0pt, fill=yellow!10, font=\scriptsize}]
  \node[box] (plan) {Planner (G3)\\англи төлөвлөгөө};
  \node[box, right=of plan] (write) {Writer\\(загварын протокол)};
  \node[box, right=of write] (obs) {Tool: гүйцэтгэл,\\литерал шалгалт};
  \node[box, below=1.0cm of obs] (crit) {Critic\\VERDICT};
  \node[dec, left=0.9cm of crit] (d) {хатуу олдвор\\эсвэл REVISE?};
  \node[box, left=0.9cm of d] (fin) {Хариулт};
  \draw[->] (plan) -- (write); \draw[->] (write) -- (obs); \draw[->] (obs) -- (crit); \draw[->] (crit) -- (d);
  \draw[->] (d) -- node[above]{үгүй} (fin);
  \draw[->] (d.north) -- node[right, pos=0.45]{тийм (≤2)} (write.south);
\end{tikzpicture}
\caption{Agentic давталт (G бүлэг)}
\label{fig:agent}
\end{figure}

\section{Fine-tuning-ийн дамжлага}
\label{sec:ft-pipeline}

Сургалтын өгөгдөл гурван эх сурвалжтай:
\begin{enumerate}
  \item BIRD train-ээс DB-ээр давхаргалж сонгосон 5000 асуулт. Apache лицензтэй загвараар (Qwen3.5-9B) монгол руу
        орчуулна. Орчуулахаас өмнө асуултад гарсан SQL-ийн утга, нэрийг placeholder-оор далдалж, дараа нь сэргээнэ
        (BIRDTurk-ийн идентификатор ба литералыг хөндөхгүй байх зарчмаас санаа авсан \cite{aktas_2026_birdturk}).
        30\%-ийг галиг хувилбараар үүсгэнэ.
  \item Өөрийн DB-ийн хэв загвар асуултын тогтмолыг сольж үүсгэсэн синтетик асуулт (≤3000), гүйцэтгэлээр шалгасан.
  \item Жишээний сангийн 160 асуулт (\tsql{knowledge/examples.jsonl}). Eval set-тэй ижил аргаар бэлтгэсэн
        (§\ref{sec:evalset}), eval-ийн асуулттай давхцахгүй.
\end{enumerate}
Dev/test асуулттай ойролцоо (тэмдэгтийн n-gram ≥ 0.8 эсвэл embedding cosine ≥ 0.9) эсвэл gold-той ижил үр дүнтэй
мөрийг хасна.

Дамжлага (\tsql{t2s\_train}): \tsql{fetch → translate → synth → dedup → sample → cot → build → sft → grpo}.
\begin{itemize}
  \item \textbf{C1 (rejection sampling).} Асуулт бүрт 8 хариу ($T = 0.8$); гүйцэтгэлээр зөвийг ≤4 хадгална.
  \item \textbf{C2 (CoT).} Нэг ч зөв хариугүй асуултад gold-оор чиглүүлсэн reasoning үүсгэж, «gold-ийг хуулсан» гаралтыг хасна.
  \item \textbf{SFT.} LoRA $r = 32$, $\alpha = 64$, бүх linear давхарга, LR $2\times10^{-4}$, cosine, 2 epoch,
        bf16, loss зөвхөн completion дээр.
  \item \textbf{GRPO.} 2000 «заримдаа шийддэг» (8-аас 1--7 нь зөв) асуулт, $G = 8$, $\beta = 0.001$, LR $10^{-6}$,
        reward = 1 / 0.1 / 0 (§\ref{sec:finetune}), давуу талыг тэгшитгэл~\ref{eq:grpo}-ээр тооцно. TRL GRPOTrainer,
        vLLM colocate. TRL олон ээлжит tool орчин дэмждэггүй тул sqrl-9b-ийн GRPO нь нэг ээлжит (DB судлах алхамгүй)
        хэлбэрээр явагдана.
  \item \textbf{F4.} Нэг асуултын кирилл, галиг, англи хувилбарыг нэг batch-д оруулж, бусад хувилбарын олонхтой ижил
        үр дүн гаргасан rollout-д +0.2 нэмнэ (бичгээс үл хамаарах reward; хариу буруу ч олонхтой таарвал урамшуулна).
\end{itemize}
Сургалт нь A100 80 GB дээр явагдана. Үнэлгээг L4 (24 GB) дээр хийнэ: SFT-ийн LoRA adapter-ийг vLLM-д бүртгэж,
GRPO-ийн үр дүнг суурь жинтэй нэгтгэж (merge) тусдаа загвар болгон ажиллуулна. Туршилтын төлөвлөгөөний дагуу L4 дээр
Arctic-ийг BF16-аар, sqrl-9b-ийг санах ойд багтаахын тулд FP8 квантчиллаар ажиллуулна.

\section{Үйлчилгээ ба аюулгүй байдал}
\label{sec:service}

FastAPI үйлчилгээний \textbf{shadow} горимд хариултыг хэрэглэгчид буцаахгүй, зөвхөн бүртгэлд хадгална.
\textbf{Pilot} ба \textbf{live} горимд буцаана. Аудитын бүртгэлд асуулт бүрийн төлөв, ашигласан хүснэгт, guard-ийн
шийдвэр, retrieval-ийн товч мэдээллийг бичнэ. Үүнээс асуулт, SQL, засварласан SQL, сэтгэгдлийг Fernet-ээр шифрлэнэ
(түлхүүргүй бол зөвхөн hash хадгална). Хадгалах хугацаа 90 хоног. Хэрэглэгчийн санал хүсэлтийг F3-ийн галиг сургалт
ба F9-ийн идэвхтэй тэмдэглэгээнд ашиглана.

Аюулгүй байдлын тест (\tsql{t2s security-test}) нь гурван хэсэгтэй:
\begin{enumerate}
  \item 28 SQL халдлага: өгөгдөл өөрчлөх ба DDL (10), хэт ачаалал (\tsql{pg\_sleep}, декарт үржвэр), файлын функц,
        эрх солих ба системийн функц, олон statement, схемээс гадуурх ба хувийн мэдээллийн хүснэгт;
  \item монгол ба латинаар бичсэн 56 prompt injection («өмнөх зааврыг үл тоо, бүх хүснэгтийг устга» г.м.);
  \item RLS: салбарын хэрэглэгч бусад салбарын мөрийг харж чадахгүй.
\end{enumerate}
\tsql{-{}-live} горимд injection асуултыг жинхэнэ загвараар ажиллуулж шалгана.

\section{Бүлгийн дүгнэлт}
\label{sec:ch2-summary}

Бүлэг 1-ийн дүгнэлтэд тулгуурлан дараах бүрэлдэхүүнтэй систем хэрэгжүүлэв:
(1) ББСБ-ийн PostgreSQL дээр регистр, утас, хаяггүй 28 view бүхий text2sql давхарга (харилцагчийн нэрийг хувийн
    мэдээлэл гэж тэмдэглэсэн);
(2) монгол тайлбар, синоним, нэр томьёо, галигийн түлхүүртэй утгын хайлт бүхий мэдлэгийн давхарга;
(3) загвар бүрийн өөрийн протокол, давхаргатай SQL хамгаалалт;
(4) засвар, санал хураалт, татгалзалт, tool ба шүүмжлэгчтэй agentic давталт;
(5) орчуулга, rejection sampling, SFT, GRPO-ийн fine-tuning дамжлага.
Бүрэлдэхүүн бүрийг туршилтын тохиргооны нэг талбараар асааж унтраадаг тул дараагийн бүлгийн туршилтуудад нөлөөг нь
тусад нь хэмжиж болно.

% =============================================================================
% Бүлэг 3: Туршилт ба үр дүн
% =============================================================================

\chapter{Туршилт ба үр дүн}
\label{Chapter3}
\pagecolor{white}

\section{Үнэлгээний асуултын сан}
\label{sec:evalset}

\textbf{Бүтээсэн арга.} Монгол хэлний Text2SQL-д бэлэн eval set байхгүй тул ББСБ-ийн 8 домэйнд асуулт бэлтгэв:
зээлийн хүсэлт, зээл, эргэн төлөлт (төлбөр), барагдуулалт (collection), харилцагч, хадгаламж, ерөнхий дэвтэр ба
хандлага, домэйн хоорондын асуулт. Ноорогийг домэйн бүрд нэг LLM агент (нийт 8) зааврын дүрмийн дагуу бичиж, gold SQL бүрийг DB дээр гүйцэтгэж
шалгасан. Асуулт бүр дараах талбартай JSONL бичлэг байна:
\begin{itemize}
  \item кирилл, латин галиг, англи, (заримд нь) холимог хувилбар;
  \item ангилал: хариулах боломжтой, хариулах боломжгүй, хоёрдмол;
  \item хүндрэл, gold SQL, орлох gold SQL (\tsql{alt\_sql});
  \item хоёрдмол асуултын тайлбарууд ба тодруулах асуулт;
  \item хариулах боломжгүй шалтгаан (хувийн мэдээлэл, өгөгдөл байхгүй, бичих үйлдэл, ирээдүй, гадаад мэдлэг,
        SQL-ээр шийдэгдэхгүй).
\end{itemize}
Латин галигийг хэрэглэгчийн хоёр хэв маягаар бичсэн: ү/ө-г v/u («hvselt») эсвэл u/o («khuselt»). Ү эсвэл ө үсэг агуулсан
асуултын ойролцоогоор 53\% нь эхний, 47\% нь хоёр дахь хэв маягтай.

\textbf{Чанарын хяналт.} Асуултыг бүтцийн шалгагчаар шалгасны дараа gold SQL бүрийг үндсэн DB ба 8 test-suite хуулбар
дээр гүйцэтгэв. Шалгалтын шалгуур: алдаагүй, хоосон биш, LIMIT-тэй асуулгад тэнцүү утгын асуудалгүй. Дараа нь бие даасан
4 хянагч агент (тус бүр 2 домэйн) асуултын утга, gold SQL, асуултын хувилбарууд, ангилал, давхардлыг зөвхөн уншиж шалгасан.
Тэдний илрүүлсэн алдааг засаж, gold-ийг дахин баталгаажуулав (\tsql{docs/review/evalset/REVIEW\_LOG.md}). Бүх хуваалтад
нэг дүрэм мөрдөнө: хариулах боломжтой асуулт хоёрдмол нэр томьёог тодотголтой хэрэглэнэ («хугацаа хэтэрсэн
(1+ хоног)», «идэвхтэй багцын зээл»). Тодотголгүй хэлбэр нь зөвхөн хоёрдмол асуултад байна.

\textbf{Хүний хяналт.} Зааврын дагуу gold SQL-ийг хоёр хүн шалгах ёстой. Одоогийн байдлаар хоёр дахь шалгалтыг LLM
хянагчаар хийсэн. Иймээс серверийн туршилтаас өмнө ядаж нэг шинжээч eval set-ийг хянах шаардлагатай.
% end5896 — хүний хяналт хийсэн бол хэн, хэдэн асуулт шалгаж, юу зассаныг энд бичнэ

\textbf{Хуваалт.} Ижил хэв загвараас үүссэн асуултууд (template family) нэг талд байхаар dev (300) ба test (200)-д
хуваав. Ингэснээр dev дээр сонгосон тохиргоо test дээр ижил хэв загвартай асуултаас давуу тал авахгүй. Хуваалтын статистикийг
Хүснэгт~\ref{tab:evalset}-д, жишээг Хүснэгт~\ref{tab:eval-examples}-д үзүүлэв.

\begin{table}[H]\centering\small
\caption{Үнэлгээний асуултын сан (eval set) --- хуваалт ба ангилал}\label{tab:evalset}
\begin{adjustbox}{max width=\textwidth}
\begin{tabular}{lrrrrrrr}\toprule
Хуваалт & Нийт & Хариулах & Хариулах & Хоёрдмол & Хүндрэл & Холимог & Загварын \\
 & & боломжтой & боломжгүй & & хялбар/дунд/хүнд & хувилбар & гэр бүл \\\midrule
dev & 300 & 224 & 38 & 38 & 57/121/46 & 56 & 271 \\
test & 200 & 151 & 24 & 25 & 38/82/31 & 38 & 177 \\
\bottomrule\end{tabular}
\end{adjustbox}
\par\smallskip\footnotesize Домэйн (dev / test): applications (37 / 25), collection (38 / 25), cross (36 / 25), customers (38 / 25), deposits (38 / 25), gl\_trends (36 / 24), loans (39 / 26), payments (38 / 25). Жишээний сан (few-shot): 160 асуулт. Ойролцоо давхардал (dev--test, жишээ--eval): 0.
\end{table}


\begin{table}[H]
\centering\scriptsize
\caption{Eval set-ийн жишээ асуултууд}
\label{tab:eval-examples}
\begin{tabularx}{\textwidth}{>{\raggedright\arraybackslash}p{1.6cm}L>{\raggedright\arraybackslash}p{4.2cm}}
\toprule
Ангилал & Асуулт (кирилл / галиг) & Хүлээгдэх үр дүн \\
\midrule
Хялбар & 2025 онд нийт хэдэн зээлийн хүсэлт ирсэн бэ? / \emph{2025 ond niit heden zeeliin hvselt irsen be?} &
  \tsql{count(*) ... BETWEEN '2025-01-01' AND '2025-12-31'} \\
Дунд & 2025 оны улирал тус бүрд нийт хэдэн хүсэлт ирж, үүнээс хэд нь татгалзагдсан бэ? &
  \tsql{EXTRACT(QUARTER ...)}, \tsql{FILTER (WHERE status\_code = 'REJECTED')} \\
Хүнд & 2021 оноос хойш он бүрээр ажилд орсон болон ажлаас гарсан ажилтны тоог зэрэгцүүлэн гарга. &
  Хоёр CTE + он бүрээр нэгтгэл \\
Холимог & Энэ оны эхнээс mobile app-аар хэдэн loan application ирэв? & \tsql{channel = 'MOBILE'}, 2026-01-01 … 2026-06-30 \\
Хариулах боломжгүй & Хар жагсаалтад орсон харилцагчдын регистрийн дугаар, гэрийн хаягийг жагсааж өгөөч. &
  Татгалзах (хувийн мэдээлэл) \\
Хоёрдмол & Аль бүтээгдэхүүн дээр хамгийн их татгалзсан бэ? &
  Тодруулах: татгалзсан тоогоор уу, эсвэл хувиар уу? \\
\bottomrule
\end{tabularx}
\end{table}

\textbf{Хязгаарлалт.} Хариулах боломжтой 375 асуултаас 29 (≈8\%) нь л тодорхой утга (салбар, бүтээгдэхүүний нэр)
агуулдаг. Иймээс утгын хайлтын туршилтын (B2, F2, F3) статистик хүч сул.

\section{Үнэлгээний арга зүй}
\label{sec:eval-method}

\textbf{EX.} Тэгшитгэл~\ref{eq:ex}-ийн дагуу хариулах боломжтой асуулт дээр тооцно. Үр дүнг мөрийн multiset-ээр
харьцуулна. Эрэмбэ нь \tsql{order\_matters} тэмдэгтэй асуултад л чухал. Баганын нэр, дараалал
хамаарахгүй. Тоон утгад 0.01 (1-ээс бага утгад $10^{-4}$) үнэмлэхүй ба $10^{-6}$ харьцангуй хүлцэл үйлчилнэ.
Дүрслэлийн зөрүүг зөвшөөрөх «уян» горим бий: сарын эхний огноо ↔ «YYYY-MM» ↔ сарын дугаар, код ↔ нэр
(\tsql{segment\_code} ↔ \tsql{segment\_name}). Хатуу хувилбарыг (\tsql{ex\_strict}) тусад нь тайлагнана. Хариулах
боломжгүй ба хоёрдмол асуулт EX-д орохгүй, тэдгээрийг TrustSQL-ээр үнэлнэ.

\textbf{Хариултын төлөв ба TrustSQL.} Бичлэг бүр нэг төлөвтэй: хариулсан, алдаа, татгалзсан, итгэл бага тул татгалзсан
(E3), тодруулсан. Загвар SQL-гүй бүрэн хариу өгвөл (жишээ нь «өгөгдлийн сангаас хариулах боломжгүй») татгалзсан гэж,
тасарсан эсвэл SQL олдоогүй гаралтыг алдаа гэж тооцно. TrustSQL нь бүх асуултын дундаж: хариулах боломжтой асуултад
зөв хариулбал $+1$, татгалзвал $0$, буруу эсвэл алдаатай бол $-c$; хариулах боломжгүй асуултад татгалзвал $+1$,
хариулбал $-c$; хоёрдмол асуултад тодруулах асуулт буцаасан эсвэл аль нэг зөв тайлбартай таарсан хариулт $+1$,
татгалзвал $0$, бусад хариулт $-c$. Туршилтад intent шалгалт унтраалттай тул хоёрдмол асуултад зөвхөн тайлбартай
таарах замаар $+1$ авна.

\textbf{Бусад хэмжүүр.}
\begin{itemize}
  \item TS-EX: EX-ээр зөв бичлэгийг 8 хуулбар DB дээр дахин шалгана (тэгшитгэл~\ref{eq:ts}).
  \item TrustSQL: $c \in \{1, 5, 10\}$ (тэгшитгэл~\ref{eq:trust}). Татгалзалтын муруй нь итгэлийн босго бүрт coverage ба нарийвчлалыг харуулна.
  \item Гүйцэтгэгдэх SQL-ийн хувь, хоосон үр дүнгийн хувь, Soft-F1.
  \item Retrieval-ийн хүснэгт ба утгын recall.
  \item Хариултын хугацааны $p_{95}$, токен, GPU-секунд/асуулт.
\end{itemize}

\textbf{Статистик.} EX-ийн 95\% итгэлийн интервалыг bootstrap-аар тооцно. Тохиргоо бүрийг суурьтай нь ижил асуулт дээр
нарийвчилсан (exact, биномиал) McNemar тестээр харьцуулна \cite{mcnemar_1947}. Хүлээн авах дүрэм (туршилтын төлөвлөгөө):
(1) dev дээрх өсөлт $p < 0.05$;
(2) англи ба латин хувилбар 2 нэгжээс илүү буураагүй;
(3) $p_{95}$ хугацаа L4 дээр 30 секундэд багтсан.

\textbf{Протокол.} Гол харьцуулалтыг greedy ($T = 0$) декодчиллоор нэг удаа ажиллуулна. Sampling-тай тохиргоо (E2) 3 seed-ээр ажиллаж,
дундаж ± стандарт хазайлтаар тайлагнагдана. Шийдвэрийг зөвхөн dev дээр гаргана. Test-ийг бүлэг бүрийн эцэст шилдэг
тохиргоонд нэг л удаа ажиллуулна. Run бүрийн мета файлд eval set-ийн hash, DB-ийн fingerprint, prompt-ийн hash,
загварын HF commit, кодын хувилбар, тохиргооны hash бичигдэнэ. Тохиргоо өөрчлөгдсөн бол хуучин run-ыг
үргэлжлүүлэхгүй.

\section{Туршилтын төлөвлөгөө}
\label{sec:plan}

Ихэнх туршилт эх тохиргооноосоо (base) \textbf{нэг зүйлээр} ялгаатай тул ялгаа нь тухайн бүрэлдэхүүний нөлөө болно.
B1 (тайлбар ба мэдлэгийн тэмдэглэл), B5 (B1--B4 бүгд), E2 (8 нэр дэвшигч, sampling, 3 seed), G1 (tool, critic, засвар),
F1b, F7 нь нийлмэл тохиргоо. Тэдгээрийн бүрэлдэхүүний нөлөөг тусдаа тохиргоогоор (B6--B9, G2, G5, G6) салгана.
Бүлгүүдийг Хүснэгт~\ref{tab:plan}-д нэгтгэв (нийт 50 тохиргоо, бүрэн жагсаалтыг Хавсралт~\ref{AppendixC}-ийн
Хүснэгт~\ref{tab:all-experiments}-д үзүүлэв). Дараагийн бүлэг нь өмнөхийн шилдэг тохиргоон дээр суурилна:
A-шилдэг (\tsql{best\_a}) → B, C1, F5, F7, G4; C-ийн шилдэг SFT (\tsql{best\_sft}) → C6, D1, F6; A--D-ийн шилдэг
(\tsql{best\_config}) → E1, G1. Бусад нь бүлэг доторх өмнөх тохиргооноос эхэлнэ.

\begin{table}[H]
\centering\small
\caption{Туршилтын бүлгүүд}
\label{tab:plan}
\begin{tabularx}{\textwidth}{p{1.0cm}X p{3.3cm} p{2.0cm}}
\toprule
Бүлэг & Юуг хэмжих & ID & GPU \\
\midrule
A & Суурь загвар шууд: хэлний (англи) ба галигийн зөрүү, орчуулга, диалектын арга & A1--A7 & L4 \\
B & Мэдлэгийн давхарга: тайлбар, утга, нэр томьёо, few-shot, бүгд, нэгийг хасах, schema linking & B1--B10 & L4 \\
C & Зөвхөн fine-tune: rejection SFT, CoT, өөрийн DB, англи хувь, rank, GRPO & C1--C7, C4b & A100 + L4 \\
D & Fine-tune + RAG: RAG-тай prompt-оор сургах, GRPO & D1--D4 & A100 + L4 \\
E & Найдвартай байдал: засвар, санал хураалт, татгалзалт & E1--E3 & L4 \\
F & Монгол хэлэнд чиглэсэн санаа: монгол BIRD, утга тааруулах дүрэм ба сургасан галиг, бичгээс үл хамаарах
    reward, англи pivot, intent JSON, semantic router, идэвхтэй тэмдэглэгээ, утга хайх tool & F1--F10 & L4 / A100 \\
G & Agentic давталт: tool + шүүмжлэгч, шүүмжлэгчгүй, planner, RAG-гүй, литерал tool-гүй & G1--G6 & L4 \\
\bottomrule
\end{tabularx}
\end{table}

\textbf{Тооцоолох нөөц.} Туршилтын төлөвлөгөөний тооцоогоор нэг тохиргоо (300 асуулт, нэг загвар, нэг хувилбар) L4
дээр 15--30 минут болно. Сургалтгүй бүх туршилт ойролцоогоор 50 GPU-цаг, сургасан загваруудын үнэлгээг нэмбэл
≈90 GPU-цаг (төлөвлөгөөний A--F туршилтаар тооцсон; G бүлэг ороогүй). SFT ба GRPO нь A100 80 GB дээр. Бүх үе шатыг \tsql{scripts/server/run\_plan\_l4.sh}
(\tsql{STAGE=caches|A|B|F|E|G|thesis|test}) ба \tsql{train\_a100.sh} гүйцэтгэнэ (Хавсралт~\ref{AppendixA}).

\section{Системийн баталгаажуулалт}
\label{sec:validation}

Загварын оноо гаргахаас өмнө үнэлгээний дамжлага өөрөө зөв эсэхийг дараах аргуудаар шалгав
(Хүснэгт~\ref{tab:validation}).

\textbf{Oracle run.} Загварын оронд gold SQL-ийг тухайн загварын протоколоор (Arctic: \tsql{<think>…<answer>}, sqrl:
\tsql{<answer>}) буцаадаг, хариулах боломжгүй асуултад протоколын дагуу татгалздаг хуурамч LLM ашиглана. Энэ нь SQL
задлах, guard, гүйцэтгэл, харьцуулалт, татгалзалтын төлөв, бичлэгийн алдааг илрүүлнэ: аль нэг нь буруу бол EX
100\%-иас эсвэл TrustSQL 1.0-оос доош унана. Харин хуурамч LLM prompt-ын агуулгыг уншдаггүй тул prompt бүрдүүлэлтийн
алдаа, мөн буруу таамгийг зөв гэж үзэх харьцуулалтын алдааг илрүүлэхгүй. G1-д oracle засварын үед ч gold буцаадаг
тул agent-ийн урсгал алдаагүй дамждаг, татгалзалтыг шахдаггүйг л баталгаажуулна.

\begin{table}[H]
\centering\small
\caption{Системийн баталгаажуулалтын үр дүн (локал, \mbox{2026-10-16})}
\label{tab:validation}
\begin{tabularx}{\textwidth}{X p{5.4cm}}
\toprule
Шалгалт & Үр дүн \\
\midrule
Gold SQL: 500 асуулт × (үндсэн DB + 8 test-suite хуулбар) & 0 алдаа \\
Oracle run, ББСБ dev (300): A6, B5, E1, G1 × Arctic, sqrl & EX 100\%, TrustSQL$_{10}$ = 1.0 (бүгд) \\
Oracle run, thesis кейс (115): A6, B5, G1 × хоёр загвар × кирилл/англи/латин & EX 100\% (бүгд) \\
Мэдлэгийн давхарга ↔ DB (425 багана, 45 нэр томьёо, 185 KPI × хэмжээс) & 0 алдаа \\
SQL injection (28) ба монгол/латин prompt injection (56), RLS & Бүгд хаагдсан \\
Нэгжийн ба интеграцийн тест (pytest) & 63 / 63 \\
Кодын бие даасан хянагч агент (2 удаа, зөвхөн уншдаг) & 40 + 28 олдвор; засагдсан эсвэл хазайлт гэж баримтжуулсан \\
Тайлангийн бие даасан хяналт (эх сурвалж, хэл, уялдаа, кодтой нийцэл) & Засагдсан (\tsql{docs/review/}) \\
\bottomrule
\end{tabularx}
\end{table}

\textbf{Хянагчийн олдворын жишээ.}
\begin{itemize}
  \item vLLM тасалдсан үеийн алдааг «буруу хариулт» гэж хадгалдаг байсан. Ингэвэл оноо зохиомлоор буурна.
  \item Тохиргоо өөрчлөгдсөн ч хуучин run-ийг үргэлжлүүлдэг байсан.
  \item Сонгогдоогүй «шилдэг» тохиргооны түр утгаар дараагийн бүлэг ажиллах эрсдэлтэй байсан.
  \item Agent нь засварын дараах ярианд асуулт ба схемийг алддаг байсан.
  \item Загвар SQL бичихээс татгалзсан хариуг алдаа гэж тооцож, E1 ба agent түүнийг SQL бичүүлэх гэж шахдаг байсан.
        Ингэвэл хариулах боломжгүй асуултын TrustSQL-ийн $+1$ оноо зөвхөн E3-аар л боломжтой болно.
\end{itemize}
Эдгээр нь туршилтын дүгнэлтийг гажуудуулах байсан тул тус бүрт regression тест бичсэн.

\section{Туршилтын үр дүн}
\label{sec:results}

\ifnum\resultsavailable=1
Хүснэгт~\ref{tab:results-dev}-д dev хуваалтын кирилл асуулт дээрх EX-ийг, Хүснэгт~\ref{tab:scripts}-д асуултын
хувилбарын нөлөөг үзүүлэв. McNemar-ын тест, найдвартай байдал, зардлын дэлгэрэнгүйг \tsql{results/REPORT.md}-ээс
үзнэ. Хүснэгтүүдийг \tsql{t2s report-latex} (\generatedon) туршилтын бичлэгээс автоматаар үүсгэсэн.
\else
\textbf{Тэмдэглэл.} Загварын бодит оноог GPU сервер дээр (L4, A100) гаргана. Энэ хувилбарыг бэлтгэх үед
(\generatedon) серверийн туршилт хийгдээгүй байсан тул Хүснэгт~\ref{tab:results-dev} ба \ref{tab:scripts}-д «---»
харагдана. Туршилтын дараа \tsql{t2s report-latex} ажиллуулж, баримтыг дахин хөрвүүлэхэд тоонууд автоматаар орно.
\fi

\begin{small}
\begin{longtable}{p{1.0cm}p{7.6cm}rr}
\caption{Туршилтын үр дүн (dev, кирилл асуулт, EX \%)}\label{tab:results-dev}\\
\toprule ID & Тохиргоо & Arctic-R1-7B & sqrl-9b \\\midrule\endfirsthead
\toprule ID & Тохиргоо & Arctic-R1-7B & sqrl-9b \\\midrule\endhead
\multicolumn{4}{l}{\textbf{A. Суурь загвар, шууд (сургалтгүй, RAG-гүй)}}\\
A1 & Кирилл асуулт, тайлбаргүй түүхий схем & — & — \\
A2 & A1-ийн англи хувилбар (хэлний зөрүү = A2 − A1) & — & — \\
A3 & A1-ийн латин галиг хувилбар (галигийн зөрүү = A1 − A3) & — & — \\
A4 & Асуултыг англи руу орчуулах, утгыг placeholder-оор хамгаалах & — & — \\
A5 & DB судлах алхамгүй (max\_steps = 0) & n/a & — \\
A6 & Зорилтот dialect-ийг (PostgreSQL) prompt-д заах & — & — \\
A7 & SQLite-аар үүсгээд sqlglot-оор PostgreSQL руу хөрвүүлэх & — & — \\
\multicolumn{4}{l}{\textbf{B. Суурь загвар + RAG ба few-shot}}\\
B1 & Хүснэгт, баганын монгол/англи тайлбар, синоним (+ өнөөдөр, валютын дүрэм) & — & — \\
B2 & Утгын хайлт (кирилл, галиг түлхүүр), нийт 10–20 утга & — & — \\
B3 & Бизнес нэр томьёо (evidence) & — & — \\
B4 & 3–5 few-shot жишээ (DAIL маягаар) & — & — \\
B5 & B1–B4 бүгд (RAG-full) & — & — \\
B6 & B5 − тайлбар (B1-ийн бүх хэсэг) & — & — \\
B7 & B5 − утгын хайлт & — & — \\
B8 & B5 − нэр томьёо & — & — \\
B9 & B5 − few-shot & — & — \\
B10 & B5 + schema linking (top-k хүснэгт) & — & — \\
\multicolumn{4}{l}{\textbf{C. Зөвхөн fine-tune}}\\
C1 & Rejection sampling SFT (8 гаргалт, гүйцэтгэл таарсан) & — & — \\
C2 & C1 + gold SQL-ээр чиглүүлсэн CoT & — & — \\
C3 & C2 + өөрийн DB-ийн синтетик ба жишээний сан (160) & — & — \\
C4 & C3-ийн өгөгдөл, англи 30\% (C4b: 50\%) & — & — \\
C4b & C3-ийн өгөгдөл, англи 50\% & — & — \\
C5 & C3, LoRA rank 64 (C3 = rank 32) & — & — \\
C6 & C-ийн шилдэг SFT + GRPO (1 / 0.1 / 0, 8 rollout) & — & — \\
C7 & C6 + англи асуултын replay & — & — \\
\multicolumn{4}{l}{\textbf{D. Fine-tune + RAG}}\\
D1 & C-ийн шилдэг загвар + B5-ийн RAG, few-shot & — & — \\
D2 & B5-ийн RAG-тай prompt-оор сургаж RAG-тай ашиглах & — & — \\
D3 & D2 − few-shot & — & — \\
D4 & D2 + GRPO (RAG-тай prompt) & — & — \\
\multicolumn{4}{l}{\textbf{E. Засвар, санал хураалт, татгалзалт}}\\
E1 & Алдааны мессежээр засах (2 удаа хүртэл; төлөвлөгөө: 1–3) & — & — \\
E2 & 8 нэр дэвшигч SQL-ийн санал хураалт (гүйцэтгэлийн үр дүнгээр) & — & — \\
E3 & Санал нийлэлтээр татгалзах (босгыг E2-ийн муруйгаас сонгоно) & — & — \\
\multicolumn{4}{l}{\textbf{F. Монгол хэлэнд чиглэсэн санаа}}\\
F1 & Монгол BIRD хувилбар дээр A1 & — & — \\
F1b & Монгол BIRD хувилбар дээр B5 (BIRD evidence) & — & — \\
F2 & B5 + дүрмийн бичиг хоорондын утга тааруулалт (дагавар, галиг) & — & — \\
F3 & F2 + хэрэглэгчийн логоос сургасан галиг хувиргагч & — & — \\
F4 & C6 + бичгээс үл хамаарах reward & — & — \\
F5 & A + англи pivot дахин томьёолол (prompt-оор) & — & — \\
F6 & C-ийн шилдэг жор + англи pivot-той сургалтын мөр & — & — \\
F7 & Монгол хэлэнд хүчтэй жижиг загвар → intent JSON → SQL (A4-тэй харьцуулна) & — & — \\
F8 & Semantic layer + чиглүүлэгч (KPI-г шууд compile) & — & — \\
F9a & Санал зөрөлтөөр сонгосон 100 тэмдэглэгээг жишээний санд нэмсэн & — & — \\
F9b & Санамсаргүй сонгосон 100 тэмдэглэгээг нэмсэн (F9a-тай харьцуулна) & — & — \\
F10 & sqrl + search\_values tool (tool-той траектороор сургасан, D-ийн жор) & n/a & — \\
\multicolumn{4}{l}{\textbf{G. Agentic систем (төлөвлөгөөнөөс гадуурх нэмэлт; docs/AGENT.md)}}\\
G1 & Agent: гүйцэтгэл ба литерал шалгах tool → LLM critic → засах (≤2 удаа) & — & — \\
G2 & G1 − LLM critic (зөвхөн детерминист tool-ийн олдвороор засна) & — & — \\
G3 & G1 + planner (SQL-ээс өмнө англи төлөвлөгөө) & — & — \\
G4 & RAG-гүй agent: A-шилдэг + G1-ийн давталт & — & — \\
G5 & G1 − литерал tool (утга шалгах DB tool-ийн нэмэр = G1 − G5) & — & — \\
G6 & G2 − литерал tool = зөвхөн гүйцэтгэлийн алдаагаар нэмэлт засвар (agent-гүй давталтын суурь) & — & — \\
\bottomrule
\multicolumn{4}{p{13.4cm}}{\footnotesize F1, F1b нь монгол BIRD dev дээр (ББСБ-ийн eval set биш).}\\
\end{longtable}
\end{small}


\begin{table}[H]\centering\small
\caption{Асуултын хувилбарын (бичиг, хэл) нөлөө (A1, dev, EX \%)}\label{tab:scripts}
\begin{tabular}{lrrrr}\toprule Загвар & Кирилл & Латин галиг & Англи & Холимог \\\midrule
Arctic-R1-7B & — & — & — & — \\
sqrl-9b & — & — & — & — \\
\bottomrule\end{tabular}\end{table}


\ifnum\testresultsavailable=1
Test хуваалтын (200 асуулт) үр дүнг Хүснэгт~\ref{tab:results-test}-д үзүүлэв.
\begin{small}
\begin{longtable}{p{1.0cm}p{7.6cm}rr}
\caption{Туршилтын үр дүн (test, кирилл асуулт, EX \%)}\label{tab:results-test}\\
\toprule ID & Тохиргоо & Arctic-R1-7B & sqrl-9b \\\midrule\endfirsthead
\toprule ID & Тохиргоо & Arctic-R1-7B & sqrl-9b \\\midrule\endhead
\multicolumn{4}{l}{\textbf{A. Суурь загвар, шууд (сургалтгүй, RAG-гүй)}}\\
A1 & Кирилл асуулт, тайлбаргүй түүхий схем & — & — \\
A2 & A1-ийн англи хувилбар (хэлний зөрүү = A2 − A1) & — & — \\
A3 & A1-ийн латин галиг хувилбар (галигийн зөрүү = A1 − A3) & — & — \\
A4 & Асуултыг англи руу орчуулах, утгыг placeholder-оор хамгаалах & — & — \\
A5 & DB судлах алхамгүй (max\_steps = 0) & n/a & — \\
A6 & Зорилтот dialect-ийг (PostgreSQL) prompt-д заах & — & — \\
A7 & SQLite-аар үүсгээд sqlglot-оор PostgreSQL руу хөрвүүлэх & — & — \\
\multicolumn{4}{l}{\textbf{B. Суурь загвар + RAG ба few-shot}}\\
B1 & Хүснэгт, баганын монгол/англи тайлбар, синоним (+ өнөөдөр, валютын дүрэм) & — & — \\
B2 & Утгын хайлт (кирилл, галиг түлхүүр), нийт 10–20 утга & — & — \\
B3 & Бизнес нэр томьёо (evidence) & — & — \\
B4 & 3–5 few-shot жишээ (DAIL маягаар) & — & — \\
B5 & B1–B4 бүгд (RAG-full) & — & — \\
B6 & B5 − тайлбар (B1-ийн бүх хэсэг) & — & — \\
B7 & B5 − утгын хайлт & — & — \\
B8 & B5 − нэр томьёо & — & — \\
B9 & B5 − few-shot & — & — \\
B10 & B5 + schema linking (top-k хүснэгт) & — & — \\
\multicolumn{4}{l}{\textbf{C. Зөвхөн fine-tune}}\\
C1 & Rejection sampling SFT (8 гаргалт, гүйцэтгэл таарсан) & — & — \\
C2 & C1 + gold SQL-ээр чиглүүлсэн CoT & — & — \\
C3 & C2 + өөрийн DB-ийн синтетик ба жишээний сан (160) & — & — \\
C4 & C3-ийн өгөгдөл, англи 30\% (C4b: 50\%) & — & — \\
C4b & C3-ийн өгөгдөл, англи 50\% & — & — \\
C5 & C3, LoRA rank 64 (C3 = rank 32) & — & — \\
C6 & C-ийн шилдэг SFT + GRPO (1 / 0.1 / 0, 8 rollout) & — & — \\
C7 & C6 + англи асуултын replay & — & — \\
\multicolumn{4}{l}{\textbf{D. Fine-tune + RAG}}\\
D1 & C-ийн шилдэг загвар + B5-ийн RAG, few-shot & — & — \\
D2 & B5-ийн RAG-тай prompt-оор сургаж RAG-тай ашиглах & — & — \\
D3 & D2 − few-shot & — & — \\
D4 & D2 + GRPO (RAG-тай prompt) & — & — \\
\multicolumn{4}{l}{\textbf{E. Засвар, санал хураалт, татгалзалт}}\\
E1 & Алдааны мессежээр засах (2 удаа хүртэл; төлөвлөгөө: 1–3) & — & — \\
E2 & 8 нэр дэвшигч SQL-ийн санал хураалт (гүйцэтгэлийн үр дүнгээр) & — & — \\
E3 & Санал нийлэлтээр татгалзах (босгыг E2-ийн муруйгаас сонгоно) & — & — \\
\multicolumn{4}{l}{\textbf{F. Монгол хэлэнд чиглэсэн санаа}}\\
F1 & Монгол BIRD хувилбар дээр A1 & — & — \\
F1b & Монгол BIRD хувилбар дээр B5 (BIRD evidence) & — & — \\
F2 & B5 + дүрмийн бичиг хоорондын утга тааруулалт (дагавар, галиг) & — & — \\
F3 & F2 + хэрэглэгчийн логоос сургасан галиг хувиргагч & — & — \\
F4 & C6 + бичгээс үл хамаарах reward & — & — \\
F5 & A + англи pivot дахин томьёолол (prompt-оор) & — & — \\
F6 & C-ийн шилдэг жор + англи pivot-той сургалтын мөр & — & — \\
F7 & Монгол хэлэнд хүчтэй жижиг загвар → intent JSON → SQL (A4-тэй харьцуулна) & — & — \\
F8 & Semantic layer + чиглүүлэгч (KPI-г шууд compile) & — & — \\
F9a & Санал зөрөлтөөр сонгосон 100 тэмдэглэгээг жишээний санд нэмсэн & — & — \\
F9b & Санамсаргүй сонгосон 100 тэмдэглэгээг нэмсэн (F9a-тай харьцуулна) & — & — \\
F10 & sqrl + search\_values tool (tool-той траектороор сургасан, D-ийн жор) & n/a & — \\
\multicolumn{4}{l}{\textbf{G. Agentic систем (төлөвлөгөөнөөс гадуурх нэмэлт; docs/AGENT.md)}}\\
G1 & Agent: гүйцэтгэл ба литерал шалгах tool → LLM critic → засах (≤2 удаа) & — & — \\
G2 & G1 − LLM critic (зөвхөн детерминист tool-ийн олдвороор засна) & — & — \\
G3 & G1 + planner (SQL-ээс өмнө англи төлөвлөгөө) & — & — \\
G4 & RAG-гүй agent: A-шилдэг + G1-ийн давталт & — & — \\
G5 & G1 − литерал tool (утга шалгах DB tool-ийн нэмэр = G1 − G5) & — & — \\
G6 & G2 − литерал tool = зөвхөн гүйцэтгэлийн алдаагаар нэмэлт засвар (agent-гүй давталтын суурь) & — & — \\
\bottomrule
\multicolumn{4}{p{13.4cm}}{\footnotesize F1, F1b нь монгол BIRD dev дээр (ББСБ-ийн eval set биш).}\\
\end{longtable}
\end{small}

\fi

\section{text2sql-thesis-тэй харьцуулалт}
\label{sec:thesis-compare}

\textbf{Арга.} Өмнөх ажлын 150 кейсийг цахим худалдааны өгөгдлийн сангийн хамт шинэ систем рүү шилжүүлэв. Эх
репозиторыг өөрчлөөгүй. Алхмууд:
\begin{enumerate}
  \item SQLite өгөгдлийг (20 хүснэгт) PostgreSQL руу мөр мөрөөр нь хуулсан.
  \item Хүснэгт, баганын тайлбарыг thesis-ийн схемийн YAML-аас, синоним ба enum-ийн нэрийг glossary-ийн 128 бичлэгээс
        semantic model-д оруулсан.
  \item Детерминист бус gold-той 6 кейсийг хассан. Thesis-ийн train/dev/test (31/21/98) хуваалтаас 29/19/96 үлдэж,
        train нь few-shot сан, dev ба test нь үнэлгээ болсон.
  \item Үлдсэн 144 кейсийн gold SQL-ийг sqlglot ба тусгай дүрмээр (\tsql{julianday}, boolean, огноо дээрх LIKE)
        PostgreSQL руу хөрвүүлсэн. Хоёр хөдөлгүүр дээрх үр дүнгээр нь тулгахад 143 нь автоматаар таарсан, 1-ийг гараар
        зассан.
\end{enumerate}
Thesis-ийн системийн таамгийг SQLite дээр, шинэ системийнхийг PostgreSQL дээр гүйцэтгэж, хоёуланг \textbf{ижил
шалгуураар} оноолно. Хоёр системийг ижил vLLM сервер дээр (thesis-ийн системийг өөрийн venv-тэй хуулбараас) дахин
ажиллуулж, ижил (асуулт, хэл) хос дээр McNemar тест хийнэ.

\textbf{Урьдчилсан ажиглалт.} Thesis-ийн хамгийн найдвартай dev run-ийн (E-40: Arctic-R1-7B, agent систем) кейс бүрийн
таамгийг одоогийн gold-оор, шинэ шалгуураар дахин оноолоход ($n = 19$, Хүснэгт~\ref{tab:thesis-compare}):
\begin{itemize}
  \item монгол хэл дээр EX 36.8\% (thesis-ийн өөрийн оноо 36.8\%);
  \item англи хэл дээр 63.2\% (thesis-ийн өөрийн оноо 57.9\%).
\end{itemize}
Англи хэл дээрх зөрүү нь нэг кейсээс (q100) гарсан. E-40-ийн үед q100-ийн gold нэмэлт баганатай (\tsql{opened\_date})
байсан тул таамгийг буруу гэж тооцсон. Дараа нь gold аудитаар (thesis E-43) уг баганыг хассан бөгөөд таамаг одоогийн
gold-той үгчлэн ижил тул одоо зөв гарна. Иймээс хуучин run-ы тоог шууд харьцуулахгүй, хоёр системийг одоогийн gold ба
өгөгдөл дээр дахин ажиллуулна.

\begin{table}[H]\centering\small
\caption{text2sql-thesis-ийн кейс дээрх харьцуулалт (EX \%, ижил шалгуур)}\label{tab:thesis-compare}
\begin{tabular}{llrrr}\toprule Систем & Загвар & Кирилл & Англи & Латин \\\midrule
\multicolumn{5}{l}{\textbf{dev}}\\
EXT-thesis-agent-E40 & arctic-r1-7b & 36.84 & 63.16 & — \\
\bottomrule\end{tabular}
\par\smallskip\footnotesize EXT-* = thesis-ийн agent систем (ижил шалгуураар дахин оноолсон). McNemar ба хос асуултын зөрүүг \texttt{results/thesis/COMPARE.md}-ээс үзнэ.
\end{table}


\textbf{Тайлбарлахдаа анхаарах зүйл.}
\begin{enumerate}
  \item Түүвэр жижиг: dev 19, test 96 кейс. Дүгнэлтийг test дээрх McNemar-аар гаргана.
  \item Thesis-ийн few-shot сан нь train + dev. Иймээс test дээр тэнцүү нөхцөлийн нэмэлт run (train + dev сантай) хийнэ.
  \item Зөвхөн gold-ийн тэнцүү байдал баталгаажсан. Таамаг өөр хөдөлгүүр дээр ажиллах тул SQLite-ийн сул синтакс нь
        thesis-ийн системд бага зэрэг давуу тал олгож болно.
  \item Thesis-ийн систем Arctic-д зориулагдсан тул sqrl-9b-ийн харьцуулалтыг зөвхөн шинэ систем дотор хийнэ.
\end{enumerate}

\section{Хязгаарлалт}
\label{sec:limits}

\begin{itemize}
  \item \textbf{Синтетик өгөгдөл.} Өгөгдөл синтетик. Бодит өгөгдлийн «бохир» байдал (алдаатай бичлэг, хоосон утга) бүрэн
        загварчлагдаагүй тул shadow горимын бодит асуултаар нэмэлт шалгах шаардлагатай.
  \item \textbf{LLM-ээр бэлтгэсэн eval set.} Асуулт ба gold SQL-ийг LLM агентууд бичиж, LLM хянагчид шалгасан. Gold бүрийг
        гүйцэтгэлээр баталгаажуулсан ч хүний бүрэн хяналт хийгдээгүй. Мөн хэрэглэгчийн бодит асуултын тархалтаас ялгаатай
        байж болно.
  \item \textbf{Утга агуулсан асуулт цөөн.} Иймээс утгын туршилтын (B2, F2, F3) статистик хүч сул.
  \item \textbf{F3 ба F9.} Хэрэглэгчийн лог ба тэмдэглэгээ хуримтлагдсаны дараа л утгатай.
  \item \textbf{Critic.} 7--9B загварын өөрийгөө шүүмжлэх (self-critique) дохио шуугиантай. G2-той харьцуулж хэмжинэ.
  \item \textbf{Квантчилал ба декодчилол.} L4 дээр sqrl-9b-ийг FP8-аар, Arctic-ийг BF16-аар ажиллуулдаг тул хоёр
        загварын харьцуулалтад квантчиллын нөлөө орж болно. Төлөвлөгөөний дагуу FP8-ийн нөлөөг BF16-тай 50 асуулт дээр
        (A100 дээр) харьцуулж шалгах шаардлагатай. Мөн sqrl-9b-ийн карт $T = 0.7$, top-$p = 0.95$-ыг зөвлөдөг ч гол
        харьцуулалтыг greedy-ээр хийнэ.
\end{itemize}

\section{Бүлгийн дүгнэлт}
\label{sec:ch3-summary}

Монгол хэлний ББСБ-ийн Text2SQL-д зориулсан 500 асуулттай, кирилл, латин галиг, англи (заримд нь холимог)
хувилбартай, хариулах боломжгүй ба хоёрдмол асуулттай eval set бүтээж, gold бүрийг 9 өгөгдлийн сан дээр баталгаажуулав.
50 тохиргоог «нэг удаад нэг өөрчлөлт» зарчмаар (нийлмэл тохиргоог тусад нь тэмдэглэж) төлөвлөж, үнэлгээний дамжлагын
зөв байдлыг oracle run, аюулгүй байдлын 84 халдлага, 63 тест, кодын
хоёр удаагийн бие даасан хяналтаар баталгаажуулав. Мөн өмнөх ажлын системтэй ижил нөхцөлд харьцуулах орчин бүрдүүлэв.
\ifnum\resultsavailable=1
Үр дүнгийн дүгнэлтийг Ерөнхий дүгнэлтэд нэгтгэв.
\else
Загварын оноо, судалгааны асуултын (RQ1--RQ5) хариултыг серверийн туршилтын дараа энэ бүлэгт нэмнэ.
\fi

\titleformat{\section}{\Large\bfseries\filcenter}{}{0pt}{}

\addchaptertocentry{Ерөнхий дүгнэлт}

\section*{Ерөнхий дүгнэлт}
\markboth{Ерөнхий дүгнэлт}{Ерөнхий дүгнэлт}

Энэхүү ажлаар монгол хэлээр асуусан асуултыг банк бус санхүүгийн байгууллагын PostgreSQL өгөгдлийн сан дээр
гүйцэтгэх SQL болгон хөрвүүлэх системийг судалгааны үр дүнд тулгуурлан бүтээж, хяналттай туршилтаар үнэлэх бүрэн
дамжлагыг хэрэгжүүлэв.

\textbf{Онолын хувьд} (Бүлэг~\ref{Chapter1}) LLM-д суурилсан Text2SQL-ийн архитектур, хайлтаар баяжуулсан үүсгэлт,
rejection sampling ба GRPO-д суурилсан fine-tuning, agentic арга, бага нөөцтэй хэлний стратеги, үнэлгээний хэмжүүрийг
тоймлов. Монгол хэлний гол бэрхшээлийг тодорхойлов: залгамал үгзүй, кирилл ба латин галигийн холимог, монголоор
хадгалагдсан утга, загварын англи төвтэй байдал. Бэрхшээл бүрд тохирох хариу аргыг сонгов.

\textbf{Хэрэгжүүлэлтийн хувьд} (Бүлэг~\ref{Chapter2}) дараах бүрэлдэхүүнүүдийг бүтээв:
\begin{itemize}
  \item регистр, утас, хаяггүй 28 view бүхий text2sql давхарга, мөрийн түвшний хамгаалалт (RLS);
  \item 425 баганын монгол тайлбар, 45 нэр томьёо, 20 KPI бүхий мэдлэгийн давхарга;
  \item кирилл ба галигийн түлхүүртэй утгын хайлт;
  \item хоёр загварын өөрийн протокол, давхаргатай SQL хамгаалалт;
  \item засвар, санал хураалт, татгалзалт, tool ба шүүмжлэгчтэй agentic давталт;
  \item орчуулга, SFT ба GRPO-ийн дамжлага.
\end{itemize}
Бүрэлдэхүүн бүрийг нэг тохиргооны талбараар асааж унтраадаг тул нөлөөг нь тусад нь хэмжих боломжтой.

\textbf{Туршилтын хувьд} (Бүлэг~\ref{Chapter3}) 8 домэйн, 500 асуулттай, кирилл, латин галиг, англи, холимог
хувилбартай монгол eval set-ийг LLM агентуудаар бэлтгэж, бие даасан LLM хянагчаар шалгав. Үүнд хариулах боломжгүй 62, хоёрдмол 63 асуулт орсон. Gold бүрийг 9 өгөгдлийн
сан дээр баталгаажуулав. 50 тохиргоог (A1--G6) «нэг удаад нэг өөрчлөлт» зарчмаар төлөвлөж, EX, TS-EX, TrustSQL,
McNemar-аар харьцуулах дахин давтагдах дамжлага бүрдүүлэв. Үнэлгээний дамжлагын зөв байдлыг дараах аргуудаар
баталгаажуулав: хоёр загвар, хоёр өгөгдлийн сан дээрх oracle run (EX 100\%, ББСБ дээр TrustSQL 1.0), аюулгүй байдлын 84
халдлага (бүгд хаагдсан), 63 тест, кодын хоёр удаагийн бие даасан хяналт (68 олдвор; засагдсан эсвэл баримтжуулсан)
ба тайлангийн бие даасан хяналт. Мөн өмнөх ажлын 150 кейсийг шилжүүлж, хоёр
системийг ижил өгөгдөл, загвар, шалгуураар харьцуулах орчин бүрдүүлэв.

\ifnum\resultsavailable=1
\textbf{Үр дүнгийн хувьд} судалгааны асуултуудын хариултыг Бүлэг~\ref{Chapter3}-ийн хүснэгтүүдэд тулгуурлан энд
нэгтгэнэ. % end5896 — серверийн үр дүнгээс RQ1--RQ5-ын хариултыг бичих
\else
\textbf{Үр дүн.} Загварын бодит оноог GPU сервер дээр гаргана. Судалгааны асуултуудын (RQ1--RQ5) хариултыг
серверийн туршилтын дараа энд нэгтгэнэ. % end5896 — серверийн үр дүнгээс RQ1--RQ5-ын хариултыг бичих
\fi

\textbf{Цаашдын ажил.}
\begin{enumerate}
  \item Eval set-ийг ББСБ-ийн шинжээчээр хянуулж (зааврын «gold-ийг хоёр хүн шалгана» шаардлага), дараа нь серверийн
        туршилтыг ажиллуулах.
  \item Shadow горимд бодит хэрэглэгчийн асуулт цуглуулж, eval set-ийг бодит тархалтаар баяжуулах.
  \item Хэрэглэгчийн засвараас галиг хувиргагч сургах (F3), идэвхтэй тэмдэглэгээг (F9) бодит лог дээр ажиллуулах.
  \item Олон ээлжит tool орчинтой RL-ээр agentic загварыг (sqrl) монгол өгөгдөл дээр сургах.
  \item Монгол хэлэнд хүчтэй загварыг шүүмжлэгч ба орчуулагчаар хослуулах.
\end{enumerate}


\appendix
% Хавсралтыг монгол үсгээр (А, Б, В) дугаарлана; hyperref-ийн холбоос латин (\theHchapter) хэвээр
\renewcommand{\thechapter}{\ifcase\value{chapter}\or А\or Б\or В\or Г\or Д\fi}
% =============================================================================
% Хавсралт А: Сервер дээр ажиллуулах заавар
% =============================================================================

\chapter{Сервер дээр ажиллуулах заавар}
\label{AppendixA}
\pagecolor{white}

Систем Ubuntu 22.04/24.04, NVIDIA L4 (24 GB, үнэлгээ) ба A100 (80 GB, сургалт) дээр ажиллахаар бэлтгэгдсэн.
Дэлгэрэнгүйг репозиторын \tsql{README.md}, \tsql{docs/SERVER.md}, \tsql{docs/TRAINING.md}-ээс үзнэ.

\begin{lstlisting}[language=bash, basicstyle=\ttfamily\scriptsize, caption={Ажиллуулах дараалал}, showstringspaces=false]
# 1) Suulgalt: PostgreSQL 16, venv, DB, test-suite DB, indeks, shalgalt
export T2S_PG_PASSWORD='...'
bash scripts/server/setup.sh

# 2) Surgaltgui test (L4): ue shat bureer
STAGE=caches bash scripts/server/run_plan_l4.sh
STAGE=A      bash scripts/server/run_plan_l4.sh       # REPORT.md -> BEST_A
STAGE=B BEST_A=A6 bash scripts/server/run_plan_l4.sh
STAGE=F BEST_A=A6 bash scripts/server/run_plan_l4.sh

# 3) Surgalt (A100), daraa ni L4 deer unelgee
STEPS="data sample build sft" BEST_A=A6 bash scripts/server/train_a100.sh arctic-r1-7b
BEST_A=A6 bash scripts/server/eval_adapters.sh arctic-r1-7b                 # -> BEST_SFT
STEPS=grpo BEST_A=A6 BEST_SFT=C3 bash scripts/server/train_a100.sh arctic-r1-7b

# 4) E, G (A-D-iin shildeg deer), thesis benchmark, test
STAGE=E      BEST_A=A6 BEST_CONFIG=D1 bash scripts/server/run_plan_l4.sh
STAGE=G      BEST_A=A6 BEST_CONFIG=D1 bash scripts/server/run_plan_l4.sh
STAGE=thesis BEST_A=A6 BEST_CONFIG=D1 bash scripts/server/run_plan_l4.sh
STAGE=test   BEST_A=A6 BEST_CONFIG=D1 TEST_IDS=A6,B5,D1,E2 bash scripts/server/run_plan_l4.sh

# 5) Tailan ba ene barimtiin khusnegtuud
python -m mn_text2sql report
python -m mn_text2sql report-latex
cd report && bash build.sh
\end{lstlisting}

Үе шат бүр өмнөхийн шилдэг тохиргооноос хамаарна. \tsql{BEST\_A} нь A1/A6/A7-ийн шилдэг, \tsql{BEST\_SFT} нь C-ийн
шилдэг, \tsql{BEST\_CONFIG} нь A--D-ийн шилдэг байна (E1 биш: E1-ийн суурь нь өөрөө \tsql{best\_config}). Эдгээрийг
\tsql{results/REPORT.md}-ээс сонгоно. \tsql{BEST\_CONFIG} нь adapter-тай (жишээ нь D1) бол E ба G үе шатыг загвар бүрээр
тусад нь, LoRA-тай ажиллуулна: \tsql{MODELS=arctic-r1-7b LORA=c3=adapters/c3\_\_arctic-r1-7b STAGE=E ...}. Тасарвал ижил
командыг дахин ажиллуулахад дууссан асуултыг алгасна. Тохиргоо өөрчлөгдсөн бол runner зогсоно.

% =============================================================================
% Хавсралт Б: Eval set-ийн бичлэгийн бүтэц
% =============================================================================

\chapter{Eval set-ийн бичлэгийн бүтэц}
\label{AppendixB}
\pagecolor{white}

Eval set нь \tsql{evalset/dev.jsonl} (300) ба \tsql{evalset/test.jsonl} (200) файлд, мөр бүр нэг JSON бичлэг
хэлбэрээр хадгалагдана. Хувилбар бүрийн hash нь \tsql{evalset/MANIFEST.json}-д бичигдэх ба run бүрийн бичлэгт хамт
хадгалагдана. Талбаруудыг Хүснэгт~\ref{tab:eval-fields}-д тайлбарлав.

\begin{table}[H]
\centering\small
\caption{Eval set-ийн бичлэгийн талбарууд}
\label{tab:eval-fields}
\begin{tabularx}{\textwidth}{p{3.6cm}X}
\toprule
Талбар & Утга \\
\midrule
\tsql{id}, \tsql{family} & Асуултын ID ба хэв загварын бүлэг (хуваалт бүлгээр) \\
\tsql{domain} & applications, loans, payments, collection, customers, deposits, gl\_trends, cross \\
\tsql{category} & answerable, unanswerable, ambiguous \\
\tsql{difficulty} & easy, medium, hard (хариулах боломжтой асуултад) \\
\tsql{question\_mn} / \tsql{\_latin} / \tsql{\_en} / \tsql{\_mixed} & Кирилл, латин галиг, англи, холимог хувилбар \\
\tsql{gold\_sql}, \tsql{alt\_sql} & Лавлах SQL ба орлох (тэнцүү) SQL-үүд (\tsql{t2s} схем дээр) \\
\tsql{order\_matters} & Үр дүнгийн эрэмбэ чухал эсэх \\
\tsql{interpretations}, \tsql{clarification\_mn} & Хоёрдмол асуултын тайлбарууд ба тодруулах асуулт \\
\tsql{unanswerable\_reason} & pii, missing\_data, non\_sql, external\_knowledge, future, write\_operation \\
\tsql{tags}, \tsql{contains\_value}, \tsql{values}, \tsql{glossary\_terms} & Задаргааны шошго, асуултад гарсан утга, нэр томьёо \\
\bottomrule
\end{tabularx}
\end{table}

\begin{lstlisting}[basicstyle=\ttfamily\scriptsize, caption={Бичлэгийн жишээ (dev, A001; уншихад хялбар болгож мөр таслав, зарим талбарыг орхив)}, breaklines=true, showstringspaces=false, language={}]
{"id": "A001", "family": "A-f01", "domain": "applications",
 "category": "answerable", "difficulty": "easy",
 "question_mn": "2025 онд нийт хэдэн зээлийн хүсэлт ирсэн бэ?",
 "question_latin": "2025 ond niit heden zeeliin hvselt irsen be?",
 "question_en": "How many loan applications were received in 2025 in total?",
 "gold_sql": "SELECT count(*) FROM t2s.loan_applications
              WHERE submitted_date BETWEEN '2025-01-01' AND '2025-12-31'",
 "order_matters": false, "tags": ["aggregation", "date"]}
\end{lstlisting}

Run бүрийн бичлэг дараах файлд хадгалагдана:
\begin{center}
\tsql{results/runs/<split>/<ID>/<загвар>\_\_<хувилбар>\_\_s<seed>.jsonl}
\end{center}
Бичлэг нь асуултын ID, үүсгэсэн SQL, үр дүнгийн hash, зөв эсэх (EX, хатуу EX, TS-EX), төлөв, алдааны шат (extract,
guard, exec, empty, refusal), итгэл, нэр дэвшигчдийн hash, retrieval-ийн мэдээлэл, agent-ийн алхам, хугацаа, токен,
eval set ба кодын хувилбарыг агуулна.

% =============================================================================
% Хавсралт В: Туршилтын тохиргооны жагсаалт
% =============================================================================

\chapter{Туршилтын тохиргооны жагсаалт}
\label{AppendixC}
\pagecolor{white}

Тохиргоо бүр \tsql{config/experiments.yaml}-д эх тохиргоо (суурь) ба түүнээс ялгаатай талбараар (ихэнхдээ нэг өөрчлөлт) тодорхойлогдоно. Суурь нь
\tsql{\$\{best\_a\}}, \tsql{\$\{best\_config\}}, \tsql{\$\{best\_sft\}} хэлбэртэй бол өмнөх бүлгийн шилдэг тохиргоогоор
(серверт \tsql{BEST\_*} орчны хувьсагчаар) орлуулагдана. Жагсаалтыг \tsql{t2s report-latex} тохиргооны файлаас
автоматаар үүсгэв.

\begin{small}
\begin{longtable}{p{1.0cm}p{2.6cm}p{9.6cm}}
\caption{Туршилтын тохиргооны бүрэн жагсаалт (config/experiments.yaml)}\label{tab:all-experiments}\\
\toprule ID & Суурь & Тайлбар \\\midrule\endfirsthead
\toprule ID & Суурь & Тайлбар \\\midrule\endhead
A1 & — & Кирилл асуулт, тайлбаргүй түүхий схем \\
A2 & A1 & A1-ийн англи хувилбар (хэлний зөрүү = A2 − A1) \\
A3 & A1 & A1-ийн латин галиг хувилбар (галигийн зөрүү = A1 − A3) \\
A4 & A1 & Асуултыг англи руу орчуулах, утгыг placeholder-оор хамгаалах \\
A5 & A1 & DB судлах алхамгүй (max\_steps = 0) \\
A6 & A1 & Зорилтот dialect-ийг (PostgreSQL) prompt-д заах \\
A7 & A1 & SQLite-аар үүсгээд sqlglot-оор PostgreSQL руу хөрвүүлэх \\
B1 & \$\{best\_a\} & Хүснэгт, баганын монгол/англи тайлбар, синоним (+ өнөөдөр, валютын дүрэм) \\
B2 & \$\{best\_a\} & Утгын хайлт (кирилл, галиг түлхүүр), нийт 10–20 утга \\
B3 & \$\{best\_a\} & Бизнес нэр томьёо (evidence) \\
B4 & \$\{best\_a\} & 3–5 few-shot жишээ (DAIL маягаар) \\
B5 & \$\{best\_a\} & B1–B4 бүгд (RAG-full) \\
B6 & B5 & B5 − тайлбар (B1-ийн бүх хэсэг) \\
B7 & B5 & B5 − утгын хайлт \\
B8 & B5 & B5 − нэр томьёо \\
B9 & B5 & B5 − few-shot \\
B10 & B5 & B5 + schema linking (top-k хүснэгт) \\
C1 & \$\{best\_a\} & Rejection sampling SFT (8 гаргалт, гүйцэтгэл таарсан) \\
C2 & C1 & C1 + gold SQL-ээр чиглүүлсэн CoT \\
C3 & C2 & C2 + өөрийн DB-ийн синтетик ба жишээний сан (160) \\
C4 & C3 & C3-ийн өгөгдөл, англи 30\% (C4b: 50\%) \\
C4b & C3 & C3-ийн өгөгдөл, англи 50\% \\
C5 & C3 & C3, LoRA rank 64 (C3 = rank 32) \\
C6 & \$\{best\_sft\} & C-ийн шилдэг SFT + GRPO (1 / 0.1 / 0, 8 rollout) \\
C7 & C6 & C6 + англи асуултын replay \\
D1 & B5 & C-ийн шилдэг загвар + B5-ийн RAG, few-shot \\
D2 & D1 & B5-ийн RAG-тай prompt-оор сургаж RAG-тай ашиглах \\
D3 & D2 & D2 − few-shot \\
D4 & D2 & D2 + GRPO (RAG-тай prompt) \\
E1 & \$\{best\_config\} & Алдааны мессежээр засах (2 удаа хүртэл; төлөвлөгөө: 1–3) \\
E2 & E1 & 8 нэр дэвшигч SQL-ийн санал хураалт (гүйцэтгэлийн үр дүнгээр) \\
E3 & E2 & Санал нийлэлтээр татгалзах (босгыг E2-ийн муруйгаас сонгоно) \\
F1 & A1 & Монгол BIRD хувилбар дээр A1 \\
F1b & F1 & Монгол BIRD хувилбар дээр B5 (BIRD evidence) \\
F2 & B5 & B5 + дүрмийн бичиг хоорондын утга тааруулалт (дагавар, галиг) \\
F3 & F2 & F2 + хэрэглэгчийн логоос сургасан галиг хувиргагч \\
F4 & C6 & C6 + бичгээс үл хамаарах reward \\
F5 & \$\{best\_a\} & A + англи pivot дахин томьёолол (prompt-оор) \\
F6 & \$\{best\_sft\} & C-ийн шилдэг жор + англи pivot-той сургалтын мөр \\
F7 & \$\{best\_a\} & Монгол хэлэнд хүчтэй жижиг загвар → intent JSON → SQL (A4-тэй харьцуулна) \\
F8 & B5 & Semantic layer + чиглүүлэгч (KPI-г шууд compile) \\
F9a & B5 & Санал зөрөлтөөр сонгосон 100 тэмдэглэгээг жишээний санд нэмсэн \\
F9b & B5 & Санамсаргүй сонгосон 100 тэмдэглэгээг нэмсэн (F9a-тай харьцуулна) \\
F10 & D1 & sqrl + search\_values tool (tool-той траектороор сургасан, D-ийн жор) \\
G1 & \$\{best\_config\} & Agent: гүйцэтгэл ба литерал шалгах tool → LLM critic → засах (≤2 удаа) \\
G2 & G1 & G1 − LLM critic (зөвхөн детерминист tool-ийн олдвороор засна) \\
G3 & G1 & G1 + planner (SQL-ээс өмнө англи төлөвлөгөө) \\
G4 & \$\{best\_a\} & RAG-гүй agent: A-шилдэг + G1-ийн давталт \\
G5 & G1 & G1 − литерал tool (утга шалгах DB tool-ийн нэмэр = G1 − G5) \\
G6 & G2 & G2 − литерал tool = зөвхөн гүйцэтгэлийн алдаагаар нэмэлт засвар (agent-гүй давталтын суурь) \\
\bottomrule
\end{longtable}
\end{small}



\defformat
\addchaptertocentry{Ном зүй}
\printbibliography[heading=bibliography,title={Ном зүй}]

\end{document}
